% ---------------------------------------------------------------------------
% Investigation-v4.tex -- merged Methods + Results, told in timeline order.
% Replaces Sections/Methods.tex and Sections/Results-and-Analysis.tex.
% All legacy \label{}s are preserved at the anchors where their content lives,
% so cross-references from the other chapters continue to resolve.
% ---------------------------------------------------------------------------

\section[The investigation]{The investigation, in the order it happened}
\label{sec:methods}
\label{sec:results}

% --- Central-claim block ----------------------------------------------------
% Each section's verdict is set off with vertical space, a quiet indent and a
% small-caps lead-in, so the spine of the argument is visible at a glance in a
% dense chapter.
\newenvironment{claim}{%
  \par\addvspace{1.4\baselineskip}%
  \begin{list}{}{\setlength{\leftmargin}{1.5em}\setlength{\rightmargin}{1.5em}%
                 \setlength{\topsep}{0pt}\setlength{\parsep}{0pt}}%
  \item\relax\noindent{\scshape The claim.}\enspace\ignorespaces
}{%
  \end{list}%
  \par\addvspace{1.4\baselineskip}%
}

This chapter merges what is conventionally split into Methods and Results, and the merger is not
cosmetic: nearly every instrument in this investigation was invented in response to the measurement
that came before it, so a methods chapter would describe tools whose entire motivation lives three
sections away, and a results chapter would quote numbers whose meaning depends on definitions the
reader met fifty pages earlier. Instead, each part below poses a question, defines the instrument
built to answer it --- with the design choices stated as the failures they avoid --- reports what
the instrument showed, and closes with what that showing forced next. Definitions therefore appear
exactly where they are first needed, and the chapter can be read straight through; each section
closes with its central claim set off from the text, so the spine of the argument is visible at a
glance. Two conventions hold throughout: every number is quoted with its interval where one exists,
and ``not established'' and ``shown to be zero'' are different statements --- the difference is
load-bearing in several places. Figure~\ref{fig:timeline} is the map.

\begin{figure}[h]
\centering
\begin{tikzpicture}[
  beat/.style={draw, rounded corners, align=center, text width=3.6cm, font=\small, inner sep=6pt},
  arr/.style={-{Latex[length=2mm]}, thick}]
\node[beat] (b1) at (0,0)    {\textbf{The ruler was wrong}\\ {\footnotesize \S\ref{sec:res-metrics}}};
\node[beat] (b2) at (4.1,0)  {\textbf{The model is drug-blind}\\ {\footnotesize \S\ref{sec:res-blind}}};
\node[beat] (b3) at (8.2,0)  {\textbf{Represented, not read}\\ {\footnotesize \S\ref{sec:methods-probe}--\ref{sec:res-mechanism}}};
\node[beat] (b4) at (12.3,0) {\textbf{Content fixes fail}\\ {\footnotesize \S\ref{sec:res-epsilon}}};
\node[beat] (b5) at (0,-2)    {\textbf{Re-encode the target}\\ {\footnotesize \S\ref{sec:methods-residual}--\ref{sec:res-encoding}}};
\node[beat] (b6) at (4.1,-2)  {\textbf{Sensitivity is not fidelity}\\ {\footnotesize \S\ref{sec:res-generalisation}}};
\node[beat] (b7) at (8.2,-2)  {\textbf{A lookup table wins}\\ {\footnotesize \S\ref{sec:res-lookup}}};
\node[beat] (b8) at (12.3,-2) {\textbf{The unseen-drug regime}\\ {\footnotesize \S\ref{sec:res-scope}}};
\draw[arr] (b1) -- (b2);
\draw[arr] (b2) -- (b3);
\draw[arr] (b3) -- (b4);
\draw[arr] (b4.south) |- (b5.east);
\draw[arr] (b5) -- (b6);
\draw[arr] (b6) -- (b7);
\draw[arr] (b7) -- (b8);
\end{tikzpicture}
\caption{The investigation in eight beats. Each is told with the instrument it required.}
\label{fig:timeline}
\end{figure}

\subsection{The atlas, the model and the rules of measurement}
\label{sec:methods-data}

The corpus is Tahoe-100M~\cite{tahoe100m}: drug-treated single-cell transcriptomes across cancer
cell lines, with plate-matched DMSO controls. Metadata supplies cell line (Cellosaurus ID), drug,
dose (\texttt{drugname\_drugconc}), plate, and a fine-grained mechanism-of-action annotation
(\texttt{moa-fine}). All analysis is restricted to the $946$ genes in the intersection of the L1000
landmark set~\cite{l1000} with Tahoe's measured genes, and expression is library-normalised to
counts-per-10k and $\log1p$-transformed.
% AUDITOR NOTE: the production normalisation is log10(1 + x) after CP10K
% (tahoe_c2s_preprocess_endcell_v2.py:214), not natural log1p. One word wrong here; anyone
% re-implementing from this text gets a different transform.

One representational pitfall is worth recording because it is silent. Tahoe's per-cell
\texttt{genes} field holds \emph{vocabulary token identifiers}, not positional indices into the gene
metadata table. Mapping them positionally scrambles gene identity quietly --- the values exceed the
metadata row count, so the error surfaces only as an out-of-bounds index at the tail. All expression
construction here maps \texttt{gene\_metadata['token\_id']} $\rightarrow$ gene symbol $\rightarrow$
panel index, and the construction is validated by a PCA cell-line probe on the resulting panel,
which reaches $0.864$ against a chance rate of $0.02$.

Three units of analysis run through everything that follows, and confusing them is the single
easiest way to misread a number. A \emph{condition} is a (drug, cell line, plate, dose) tuple ---
the level at which responses are scored. The physical treatment assignment is the \emph{sample},
one treated well, and one well holds roughly $49$ cell lines dosed together, so a condition is an
observation nested inside an assignment, not an assignment. That distinction is cosmetic here and
decisive in \S\ref{sec:methods-holdout}. Plates are numbered within cell line rather than globally,
so the batch unit is the pair (cell line, plate) and never the plate label alone. Dose is carried
as a molar concentration resolved from \texttt{drugname\_drugconc}, not as a unitless float: an
earlier pipeline kept only the numeric part, which collides $1\,\mu$M with $1\,$nM, and a direct
audit of all $289$ treatment samples found $289$ usable molar doses with zero unit collisions ---
every drug in this atlas states its doses in a single unit, so the defect was a demonstrated
possibility that did not materialise. Combination samples are dropped rather than analysed as their
first component.

Table~\ref{tab:scale} states the quantities every $n$ in this chapter should be located against.
Two populations are used and they are not interchangeable. The broad characterisation of the task
(\S\ref{sec:res-blind}) streams a wide sample of the atlas; the residual arm (\S\ref{sec:res-encoding}
onward) works from a smaller cached subset built once and reused, so that every arm sees identical
cells.

\begin{table}[h]
\centering
\begin{tabular}{lr}
\toprule
\multicolumn{2}{l}{\textbf{Task characterisation} (streamed sample)} \\
\midrule
conditions characterised & 6,628 \\
\quad of which statistically active against DMSO & 3,647 / 5,014 tested (72.7\%) \\
\quad identifiable in principle (ceiling $\geq 0.8$) & 3,064 (46.2\%) \\
\quad with a plate-mate closer than their own replicate & 78.7\% \\
median cells per condition & 44 \\
median \#DEGs per condition (BH $q<0.05$) & 0 \\
median SNR (effect $\div$ replicate noise) & 0.75 \\
drugs with a real dose series & 218 / 351 (62.1\%) \\
drugs seen in $\geq 3$ cell lines & 227 \\
\midrule
\multicolumn{2}{l}{\textbf{Residual-frame cache} (rebuilt targets, \S\ref{sec:methods-residual})} \\
\midrule
cells & 620,564 (600,000 treated, 20,564 control) \\
(cell line, plate) groups & 198 \\
drugs / cell lines & 106 / 50 \\
conditions with $\geq40$ treated and $\geq20$ control cells & 6,617 \\
\quad by split after inventory filters (train / held-out wells / held-out drugs) & 4,999 / 900 / 711 \\
\quad retained by the data-resolved reliability line & 2,523 of 4,999 train (50.5\%) \\
training examples / well-held-out validation examples & 147,710 / 3,600 \\
\midrule
\multicolumn{2}{l}{\textbf{Evaluation sets}} \\
\midrule
tier 1 (seen conditions) / tier 2 (unseen drugs) / tier 3 (unseen combinations) & 4,325 / 2,188 / 99 \\
conditions scored per tier in the generation evaluations & 400 \\
tier-2 drugs in the expression-frame $\NIR$ benchmark (within plate) & 606 \\
conditions scored in the residual-frame evaluation & 570 (41 cell lines, 87 drugs) \\
\bottomrule
\end{tabular}
\caption{The populations behind every $n$ quoted in this chapter. Sources:
\texttt{drug\_biology\_atlas.json}, \texttt{eval\_endcell\_cpu.json},
\texttt{nir\_sameplate.json}, the rebuilt target report and the residual-frame evaluation. Note
that the 72.7\% active figure is over the 5,014 conditions with enough cells for the permutation
test, not over all 6,628.}
\label{tab:scale}
\end{table}

Three evaluation tiers, defined by what the training set contained, are used for the broad
characterisation. \emph{Tier 1} holds conditions whose (drug, cell line) pairing was seen in
training, so it measures fit rather than generalisation. \emph{Tier 2} holds conditions whose
\emph{drug} never appeared, which is Leave-Drugs-Out in the vocabulary of
Section~\ref{sec:lit-dreval}. \emph{Tier 3} holds pairings whose drug and cell line were each seen
but never together --- Leave-Pairs-Out. Tier 3 is small ($99$ conditions) and is reported only
where it is adequately powered; the residual arm builds its own, larger holdout, and the unit that
holdout is keyed at turns out to be one of the most consequential choices in the chapter
(\S\ref{sec:methods-holdout}).

The model side is equally compact. A cell is written as a \emph{cell sentence}: its expressed panel
genes in descending order of expression, terminated by \ENDCELL. The prompt supplies cell line,
drug, dose, mechanism, and the untreated control cell's own sentence; the response is the treated
cell's sentence. Two tokens, \ENDCELL{} and --- for the residual arm --- \DOWN{}, are registered as
atomic special tokens; left unregistered, \DOWN{} splits into sub-words and the up/down boundary
ceases to be a clean symbol. The backbone is \texttt{C2S-Scale-Pythia-1b-pt}~\cite{c2sscale}, a
C2S-Scale checkpoint built on the Pythia-1b decoder~\cite{pythia}, cold-started for every arm from
the same base checkpoint with identical hyperparameters (1 epoch, learning rate $10^{-5}$, gradient
accumulation 16, maximum length 8192, bf16), so that in each experiment the training target or
objective is the only variable. What that checkpoint has already seen bears directly on how the
held-out splits should be read: it is pretrained on over 50 million human and mouse cells drawn
from 825 single-cell datasets in CELLxGENE and the Human Cell Atlas --- an \emph{atlas} corpus,
carrying no chemical or genetic perturbation data, and the released model card lists cell and
tissue prediction, conditioned generation and gene-set conversion among its capabilities but not
perturbation prediction. The base model therefore reaches our fine-tuning having never seen a
drug-treated cell, which is what allows the unseen-drug split to mean what it claims: a drug held
out here is unseen by the whole pipeline, not merely by its final stage.

Two rules of measurement apply to every result below, and are collected here so the later sections
can quote numbers without re-justifying them. The first concerns uncertainty. Conditions are not
independent, and they are not independent in two crossed ways at once: two conditions from the same
treated well share a physical assignment and a batch, and two conditions from the same cell line
share the biology. Treating conditions as independent gives intervals that are too narrow --- here,
about $1.3$ to $1.5$ times too narrow --- and neither grouping alone captures the dependence. The
tempting shortcut is actively misleading: clustering on the assignment unit alone uses the
\emph{larger} cluster count ($289$ wells against $\approx 50$ lines) and would \emph{tighten}
every interval while sounding more rigorous. All intervals in this chapter are therefore two-way
cluster-robust (Cameron--Gelbach--Miller),
\begin{equation}
  V \;=\; V_{\textrm{line}} \;+\; V_{\textrm{well}} \;-\; V_{\textrm{intersection}},
\end{equation}
where the intersection of the two groupings is the single condition, so $V_{\textrm{intersection}}$
is the ordinary independent-sampling variance; when the estimator goes non-positive in a finite
sample the wider one-way variance is used and the branch recorded. Because a cluster-robust
variance is asymptotic in the number of \emph{clusters}, not observations, the reference
distribution is Student's $t$ with degrees of freedom one less than the smaller cluster count ---
which matters where an arm spans only $25$ wells, and moves a borderline verdict. For statistics
computed over \emph{pairs} of conditions, such as the transfer coefficient, neither grouping
partitions the observations --- a pair belongs to two nodes at once --- so those intervals use the
Fafchamps--Gubert dyadic estimator, under which two pairs are dependent whenever they share either
member. This is the structure DrEval~\cite{dreval} identifies in this field more broadly, citing
Hurlbert's pseudoreplication~\cite{hurlbert}: hundreds of thousands of cells resolving to a few
hundred conditions over $50$ cell lines.

The second rule concerns generation validity. Every generation-based evaluation reports the
\ENDCELL{}/\DOWN{} completion rate, the fraction of emitted symbols that are valid panel genes, and
the duplicate rate, before any score is quoted. This is not hygiene: a $600$-token generation cap
silently truncated $26\%$ of generations in an early run and \emph{halved} the measured effect.
Every contrast is additionally generated with a fixed per-(condition, arm) seed, so the two arms of
a paired comparison are not drawn from different points of one global stream.

\subsection{The metric is the first result}
\label{sec:methods-metrics}
\label{sec:res-metrics}

The investigation begins with the scoreboard, because every later claim depends on it. The field's
default metric is a correlation between predicted and true \emph{change from control}, restricted
to the genes that actually move. Writing $x$ for the control profile, $y$ for the treated truth and
$\hat{y}$ for the prediction, let
\begin{equation}
  S_K \;=\; \operatorname*{arg\,top-}K_{\,g} \; \lvert y_g - x_g \rvert
\end{equation}
be the $K$ genes with the largest true change. Then
\begin{equation}
  \DEdr \;=\; \operatorname{corr}\!\big( \left(\hat{y} - x\right)\big|_{S_K}, \;
                                          \left(y - x\right)\big|_{S_K} \big).
\end{equation}
We report $K=50$ with Pearson correlation throughout; $K \in \{20, 100, 200\}$ and the Spearman
variant were computed and move the numbers by less than $0.02$ except at $K=200$, where the
top-$K$ set starts to include genes that do not move.

What does this metric reward? The measurement is a predictor that contains nothing: every gene
placed at the middle rank, carrying no information about drug, cell or biology. It scores
$\DEdr(K{=}50) = 0.9999$ [$0.99991$, $0.99992$] --- above the two-real-cell noise ceiling ($0.76$)
and above the model ($0.73$). The ordering is inverted: less information yields a higher score,
which is the signature of an artefact rather than of skill. The same predictor's partial-$\DEdr$
and $\paneltau$ are \emph{undefined} rather than merely poor --- its shift from control is
identically zero, so there is nothing left to correlate once the control is regressed out.

\begin{table}[h]
\centering
\begin{tabular}{lrrr}
\toprule
Predictor & Information & $\DEdr$ (K50) & partial-$\DEdr$ \\
\midrule
\texttt{revert\_center} (all genes $\rightarrow P/2$) & none & \textbf{0.9999} & undefined \\
\texttt{revert\_mean} (predict mean control) & none, no fit & 0.961 & 0.25 \\
ridge linear (control $\rightarrow$ shift) & control & 0.947 & 0.28 \\
noise ceiling (two real cells) & --- & 0.76 & --- \\
model (1B LLM) & control + drug & 0.73 & not measured \\
\bottomrule
\end{tabular}
\caption{$\DEdr$ is saturated by a control artefact. Regressing the control out of the two fitted
baselines leaves $\approx 0.26$ of real but \emph{drug-agnostic} skill, and a baseline that performs
no fitting at all reaches the same figure. The model's own partial-$\DEdr$ was never computed --- it
requires a GPU pass that was not run --- so the cell is marked rather than filled with the expected
value; the model's $\paneltau$, which \emph{is} measured, carries the same conclusion below.}
\label{tab:de-artifact}
\end{table}

The cause is visible in the definition, and three properties of it, each innocent in isolation,
compose into the failure. First, $S_K$ is selected using the \emph{truth}, so every predictor is
scored on the genes the condition actually moved --- the exam questions are chosen from the answer
key. Second, treated cells in this atlas regress toward a shared response, so a prediction that
simply reverts toward the population centre correlates with $y - x$ \emph{by construction}: no drug
knowledge is required to score well, only knowledge of where the centre is. Third, the fix is
available in the same breath --- \textbf{partial-}$\DEdr$ regresses the control profile out of both
arguments before correlating, which removes the reversion route and is therefore the honest version
of the same quantity. Removing the control collapses the raw $0.95$ to $\approx 0.26$ for the two
baselines that were scored --- and that residue is matched by a baseline that performs no fitting at
all. The model's own partial-$\DEdr$ is not measured, so the comparison is anchored instead to
$\paneltau$ --- Kendall's $\tau$ between predicted and true rank over the full $946$-gene panel,
which selects no gene subset --- measured for every arm and telling the same story:
\texttt{revert\_mean} $0.268$, ridge $0.265$, model $0.257$ (tier 2, \texttt{worst} convention).
The response decomposes into control reversion (an artefact), a generic drug program
($\approx 0.26$, drug-agnostic, trivially matched), and a drug-specific component that no predictor
captures.

A cell sentence lists only expressed genes, so a predicted profile leaves some panel genes
unranked, and any rank-based metric needs a convention for the absent tail. Two were used
throughout and reported side by side: \texttt{worst} assigns every absent gene the last rank $P$,
and \texttt{francesca} assigns them a fixed middle rank $P/2$. The choice is immaterial --- spike-in
$\DEdr$ reads $0.952$ against $0.954$ under the two, and no baseline or ceiling in this chapter
moves by more than $0.01$ --- so all figures are quoted under \texttt{worst} and the convention is
not mentioned again.

The metric this investigation relies on instead is the Normalised Inverse Rank, introduced by Wu
et al.\ in PerturBench~\cite{perturbench} (there called the ``transposed-rank'' metric) and
identified by Miller et al.~\cite{miller} as one of the few consistently calibrated metrics across
14 Perturb-seq datasets. For a prediction $\hat{y}$ of condition $i$, with truths $\{y_j\}$ over
the comparison set,
\begin{equation}
  \NIR(\hat{y}, i) \;=\; \frac{1}{|J|}\sum_{j \in J}
    \Big[\mathbf{1}\{\sigma(\hat{y}, y_i) > \sigma(\hat{y}, y_j)\}
       + \tfrac{1}{2}\,\mathbf{1}\{\sigma(\hat{y}, y_i) = \sigma(\hat{y}, y_j)\}\Big],
\end{equation}
the fraction of \emph{other} conditions whose truth the prediction resembles less than its own,
with ties counted at one half. Two properties do work that no other candidate does. Chance is
$0.50$ \emph{by construction}, not empirically: ``this model is at chance'' is a statement, not a
comparison against a fitted null, and that is what makes the negative results in this chapter
quotable without a control arm. And because the generic response is present in every condition's
truth in equal measure, it cancels in the comparison --- $\NIR$ is a \emph{discrimination} metric
and only the drug-specific difference decides it, which is precisely the component the reversion
artefact exploits. The tie term is not cosmetic: in residual space
(\S\ref{sec:methods-residual}) a predictor of ``the average drug response'' is exactly the zero
vector, every similarity against it is zero, and a strict inequality scores it $0.000$ rather than
the $0.500$ it deserves.

Our instruments compute $\NIR$ under two similarity functions and store both: \texttt{nir\_expr},
Euclidean distance in expression space, and \texttt{nir\_rank}, the same construction over rank
vectors. \textbf{Every $\NIR$ figure in this chapter is \texttt{nir\_expr}}, and the choice is
load-bearing rather than incidental: under \texttt{nir\_rank} the tier-2 aggregate places
\emph{every} arm below chance, including the noise ceiling itself ($0.380$, against the model's
$0.313$ and a drug-agnostic linear map's $0.285$). A metric on which a real replicate of a
condition fails to be retrieved above chance is not measuring what its name suggests, and it is not
the metric the calibration analysis below endorsed --- the $\DRF$ result is computed on the
expression-space variant. We report the calibrated one and name it here so the stored artifacts
cannot be misread.

% fig-nir.tex -- NIR as a discrimination score. Goes with the metric definition.
% No data dependency.
\begin{figure}[htbp]
\centering
\begin{tikzpicture}[
  x=1cm, y=1cm, font=\small,
  lbl/.style={font=\scriptsize, text=black!70},
  strong/.style={font=\scriptsize, text=black!90},
  own/.style={circle, fill=black!85, inner sep=2pt},
  other/.style={circle, draw=black!55, fill=white, inner sep=1.9pt},
  ar/.style={-{Stealth[length=4.5pt]}, black!65},
]

\node[lbl, anchor=south] at (2.05,0.44) {other drugs in the same comparison set};
\node[strong, anchor=south] at (6.30,0.44) {own truth};

\fill[black!10, rounded corners=1pt] (0.55,0.00) rectangle (6.30,0.34);
\foreach \x in {0.95, 1.75, 2.60, 3.55, 4.45, 5.45} {\node[other] at (\x,0.17) {};}
\node[own] at (6.30,0.17) {};
\node[other] at (7.15,0.17) {};
\node[other] at (7.85,0.17) {};
\node[anchor=west, font=\small] at (8.35,0.17) {$\NIR = \tfrac{6}{8}$};

\draw[ar] (0.45,-0.26) -- (8.15,-0.26);
\node[lbl, anchor=north] at (4.30,-0.38)
  {similarity of the prediction to each condition's truth $\longrightarrow$};

\draw[decorate, decoration={brace, mirror, amplitude=3.5pt}, black!55]
  (0.55,-0.80) -- (6.30,-0.80);
\node[strong, anchor=north] at (3.40,-0.96) {the shaded fraction \emph{is} the score};

\node[lbl, anchor=west] at (0.45,-1.50) {own truth ranked first $\Rightarrow \NIR \to 1$};
\node[lbl, anchor=west] at (0.45,-1.84) {own truth ranked at random $\Rightarrow \NIR = 0.5$};
\node[strong, anchor=west] at (0.45,-2.18)
  {every similarity \emph{identical} $\Rightarrow \NIR = 0.5$ by the half-tie rule, not $0$};

\end{tikzpicture}
\caption[NIR is a discrimination score]{\textbf{$\NIR$ asks whether a prediction resembles its own
condition's truth more than it resembles the truths of the other drugs in the same comparison set.}
The generic response --- the stress and cell-cycle program every compound triggers --- is present in
every truth in equal measure, so it cancels in the comparison and only the drug-specific difference
decides. Two consequences follow, and both matter downstream. Chance at $0.50$ is a statement about
the \emph{comparison set}, not about the model: conditions are grouped by (cell line, plate) in the
expression frame and by cell line in the residual frame, and the two are not interchangeable. And a
predictor whose similarities are all identical --- the zero vector, which is what ``predict the
average drug response'' becomes in residual space --- must score $0.50$ rather than $0$, which is why
ties are counted at one half rather than dropped.}
\label{fig:nir}
\end{figure}


Choosing $\NIR$ over its rivals is itself a measurement, made with the framework of Miller
et al.~\cite{miller}, which grades each \emph{metric} rather than each model by asking how well it
separates a positive from a negative control:
\begin{equation}
  \DRF(m) \;=\; \frac{m(\textrm{pos}) - m(\textrm{neg})}{m(\textrm{perfect}) - m(\textrm{neg})},
\end{equation}
with $\textrm{pos}$ an interpolated-duplicate noise ceiling --- their positive control, which
interpolates a technical duplicate toward the mean baseline with per-gene weights given by
differential-expression significance --- and $\textrm{neg}$ a drug-agnostic leave-one-out mean. A
negative \DRF{} means the metric scores the uninformative baseline \emph{above} the noise ceiling:
it is inverted, and rewards the generic gene ordering rather than the drug. Three hardening choices
matter for how the result can be read. Tied expression values receive \emph{average} ranks, because
breaking ties by array index makes a tied gene's rank depend on its position in the panel. The
\DRF{} ratio is recomputed \emph{whole} inside every bootstrap resample, because its denominator is
estimated from the same data as its numerator and holding it fixed understates the interval. And
the family of five metrics is controlled simultaneously with Holm's procedure, because the claim
``only $\NIR$ is calibrated'' is a claim about all five at once.

Under the hardened protocol, $\NIR$ is the only calibrated metric: $\DRF =
\ci{+0.635}{+0.618}{+0.652}$, Holm $p = 0.0025$. The four rivals remain strongly negative and none
approaches significance (Holm-adjusted $p = 1.0000$ for all four); in the earlier unhardened
calibration they read $-0.081$ (\texttt{weighted\_r2}), $-0.231$ ($\paneltau$), $-0.415$
(\texttt{spearman\_expr}) and $-0.448$ ($\DEdr$). The inversion is a property of the data, not a
leak: a leave-one-out correction barely moves the numbers, since real cell lines carry 20--204
drugs, and it survives holding the plate constant.

This replicates Miller et al.\ --- and then diverges from them. Their central methodological claim
is confirmed emphatically: metric choice dominates, the rank-based metric is the one that survives,
and the correlation-based metrics they identify as poorly calibrated are exactly the ones that fail
here. The agreement extends to the specific metrics --- they find $\NIR$ consistently calibrated
across 14 Perturb-seq datasets, and we find it is the only calibrated metric of the five we tested
on drug perturbation. Our \texttt{weighted\_r2} result is also consistent with their finding that
weighted metrics behave better than unweighted ones. The divergence comes at the next step. Miller
et al.\ conclude that under calibrated metrics deep models outperform uninformative baselines;
\S\ref{sec:res-blind} evaluates on the calibrated metric and finds the model at chance. Two
differences in setting are candidates for the discrepancy --- genetic versus drug perturbation, and
a discriminative regression model versus an autoregressive generative one --- and
\S\ref{sec:res-lookup} supplies a third and more structural possibility. Adopting their framework
is also deliberate in a second sense: evaluating our model on the metric \emph{their} analysis
endorses means the drug-blindness result cannot be dismissed on the grounds their own paper raises.

\begin{figure}[htbp]
\centering
\includegraphics[width=\textwidth]{fig-calibration}
\caption[The standard metric is exploitable]{\textbf{Under the standard metric, less information
yields a higher score.} (a) A predictor that places every gene at the middle rank --- carrying
nothing about drug, cell or biology --- scores $\DEdr = 0.9999$, above the two-real-cell noise
ceiling and above both the model and its own scramble. The inverted ordering is the signature of an
artefact rather than of skill, and its cause is control regression-to-the-mean. Tier 2, $n=400$
conditions. No intervals are drawn: the stored ones are normal approximations rather than the
cluster-robust intervals used everywhere else in this chapter, and mixing the two conventions
on one axis would misrepresent both. (b) Grading the \emph{metrics} rather than the models. The
displacement drawn is the Dynamic Range Fraction, on a scale where $0$ is the drug-agnostic baseline
and $1$ a perfect predictor, so a negative value means the metric ranks an uninformative baseline
\emph{above} the noise ceiling. Only $\NIR$ points forward; the hardened estimate is
\ci{+0.635}{+0.618}{+0.652}. $655$ drugs over $61$ cell-line$\times$plate
groups, within plate. \texttt{weighted\_r2} carries no directional claim --- its sign tracks cell
count, and it is approximately neutral.}
\label{fig:calibration}
\end{figure}

The metric can discriminate even if the model cannot, and the instrument that shows it is a
spike-in forced choice: within one plate, a reference pseudobulk drawn from drug A's cells must be
scored more similar to a disjoint sample of A's own cells than to a sample of drug B's, over many
resamples, with chance at $0.50$; contaminating B's population with a growing fraction of A's cells
titrates the difficulty, so the curve of accuracy against contamination characterises each metric's
sensitivity. At pseudobulk and zero contamination the separation is essentially perfect
($\paneltau = 1.000$, $\DEdr = 0.995$, top-$n$ $\tau = 0.987$, plate held constant). The instrument
is not the bottleneck.

\begin{claim}
The field's standard score is saturated by a predictor carrying no information ($\DEdr = 0.9999$);
of the candidates audited, only $\NIR$ is calibrated ($\DRF = \ci{+0.635}{+0.618}{+0.652}$); and on
$\NIR$ the instrument can discriminate. Every drug-use claim in the rest of this chapter is
therefore made on the calibrated metric.
\end{claim}

\subsection{The model is drug-blind}
\label{sec:methods-controls}
\label{sec:res-blind}

With a ruler that cannot be gamed in hand, the investigation can finally ask the question it exists
to answer --- does the model use the drug at all? --- and the answer is uniform across four
instruments that fail in different ways, which is the only reason to report all four. Each
instrument exists because a specific objection would otherwise be fatal, and each is stated with
the failure it closes.

Drug and plate are partially confounded by Tahoe's design: each drug is assayed on its own
plate(s), with roughly 4--12 drugs per plate and plate-matched controls. Discrimination grouped by
cell line alone therefore lets the \emph{batch} signature stand in for drug identity. Measured
directly, by scoring the same tier-2 conditions with and without the plate held constant:

\begin{table}[h]
\centering
\begin{tabular}{lrr}
\toprule
Predictor (drug-agnostic) & cross-plate & within plate \\
\midrule
mean baseline & 0.399 & \textbf{0.161} \\
control-copy & 0.551 & 0.510 \\
\midrule
noise ceiling & 0.625 & 0.575 \\
fraction of drugs at chance & 43\% & 49\% \\
\bottomrule
\end{tabular}
\caption{Expression-frame $\NIR$. A drug-agnostic mean baseline more than doubles its score when the
plate is free to vary, which is what plate leakage looks like. The ceiling falls once the plate is
held constant because the batch signature was inflating it too.}
\label{tab:plateleak}
\end{table}

The mean baseline is the clearest demonstration: it contains no drug information whatsoever, and it
scores $0.399$ when the comparison set may span plates against $0.161$ when it may not. Every
discrimination instrument therefore gained a \texttt{--same\_plate\_only} mode and all discrimination
results in this chapter were re-run within plate. Conclusions are unchanged; some magnitudes shrink.

The central manipulation is the \emph{scramble ablation}: the prompt is left byte-identical except
that the drug name and mechanism are replaced by another drug's, while the control cell and the
scoring truth are unchanged. Because both arms share the same control cell, control- and
batch-derived leakage cancels in the \emph{difference} $\gap = \NIR_{\textrm{model}} -
\NIR_{\textrm{scramble}}$. This is the only manipulation that isolates drug use, and every
drug-use claim in this chapter is anchored to it rather than to a comparison against a
drug-agnostic predictor.

Two constraints on the swap partner are required for a gap measured this
way to mean what it says, and both exist because the unconstrained pool fails in a specific way.
The partner must be a \emph{different drug}: allowing the same drug at another dose or plate
substitutes the drug for itself, an unchanged output is then correct rather than blind, and the
contamination concentrates in the most similar stratum because a drug's own other doses are usually
its nearest neighbours. And the partner must be on the \emph{same plate}: residuals retain plate
structure, so a cross-plate partner differs from the target in batch as well as in drug, and the
gap absorbs both. A scramble that swaps in a drug with a \emph{similar} true response is also a
weak test, and this is a real risk here because same-mechanism drugs are barely more similar than
random pairs (within-MoA distance $2.322$ versus between-MoA $2.377$, ratio $0.977$), so
``different mechanism'' does not imply ``different response''. Two remedies are used: a
\emph{swap-distance sweep} --- two hundred (drug, swap) pairs over twenty cell lines on tier 2,
binned by quantile of the response distance $\lVert y_A - y_B \rVert$ between the real and
swapped-in drug --- which turns the single test into a curve of $\gap$ against that distance; and a
\emph{stratified scramble} with three defined partners per condition chosen by residual cosine
within the cell line --- \texttt{near} (most similar), \texttt{orth} ($\cos\approx0$) and
\texttt{opposite} (most anti-correlated). The strata make an assumption visible that a single comparator hides: the gap has
two parts, how much naming the truth helps and how much naming a lie hurts, and only the first is
wanted. A usable null is therefore a stratum whose own score sits at chance --- which,
\S\ref{sec:res-generalisation} shows, is true of exactly one of the three, and not the one an
evaluation would naturally default to.

% fig-scramble.tex -- the scramble ablation and why its difference is leak-immune.
% No data dependency.
% Layout note: the prompt boxes are ~2.2cm tall once the shaded spans are set, which is far more
% than a four-line block looks like it should need. Every band below them is placed against that
% measured height, not against an estimate.
\begin{figure}[htbp]
\centering
\begin{tikzpicture}[
  x=1cm, y=1cm, font=\small,
  lbl/.style={font=\scriptsize, text=black!70},
  strong/.style={font=\scriptsize, text=black!90},
  promptbox/.style={draw=black!35, rounded corners=2pt, inner sep=5pt, fill=black!2},
  ar/.style={-{Stealth[length=4.5pt]}, black!65},
]
\newcommand{\spanhl}[1]{\colorbox{black!14}{#1}}

\node[lbl, anchor=west] at (0.15,0.00) {real prompt};
\node[promptbox, anchor=north west] (p1) at (0.15,-0.26) {%
  \renewcommand{\arraystretch}{0.9}%
  \begin{tabular}{@{}l@{~~}l@{}}
  \ttfamily\scriptsize Cell line:    & \ttfamily\scriptsize MCF7 \\
  \ttfamily\scriptsize Drug:         & \ttfamily\scriptsize \spanhl{Lapatinib} \\
  \ttfamily\scriptsize Mechanism:    & \ttfamily\scriptsize \spanhl{EGFR inhibitor} \\
  \ttfamily\scriptsize Control cell: & \ttfamily\scriptsize ACTB GAPDH \ldots \\
  \end{tabular}};

\node[lbl, anchor=west] at (0.15,-2.86) {scrambled prompt --- only the two shaded spans differ};
\node[promptbox, anchor=north west] (p2) at (0.15,-3.12) {%
  \renewcommand{\arraystretch}{0.9}%
  \begin{tabular}{@{}l@{~~}l@{}}
  \ttfamily\scriptsize Cell line:    & \ttfamily\scriptsize MCF7 \\
  \ttfamily\scriptsize Drug:         & \ttfamily\scriptsize \spanhl{Vorinostat} \\
  \ttfamily\scriptsize Mechanism:    & \ttfamily\scriptsize \spanhl{HDAC inhibitor} \\
  \ttfamily\scriptsize Control cell: & \ttfamily\scriptsize ACTB GAPDH \ldots \\
  \end{tabular}};

\node[draw=black!35, rounded corners=2pt, inner sep=6pt, align=center, fill=black!3]
  (truth) at (9.30,-2.55) {truth for\\\textbf{Lapatinib}\\{\scriptsize drawn once}};

\draw[ar] (p1.east) .. controls (6.6,-1.35) and (7.6,-1.85) .. (truth.north west);
\draw[ar] (p2.east) .. controls (6.6,-4.20) and (7.6,-3.30) .. (truth.south west);

\node[anchor=north, font=\small] at (9.30,-3.70)
  {$\gap = \NIR_{\text{real}} - \NIR_{\text{scrambled}}$};

\draw[ar, dashed, black!45] (1.20,-6.10) -- (8.10,-6.10);
\node[strong, anchor=south, align=center] at (4.65,-6.02)
  {the shared control cell enters \emph{both} arms, so control- and batch-derived};
\node[strong, anchor=north, align=center] at (4.65,-6.20)
  {leakage cancels in the difference};

\end{tikzpicture}
\caption[The scramble ablation]{\textbf{The scramble replaces the drug name and its mechanism and
nothing else.} Both arms are conditioned on the same control cell and scored against the same truth,
so any leakage through the control cell or the plate enters both arms in equal measure and cancels in
the difference $\gap$. This is why $\gap$, rather than a comparison against a drug-agnostic
predictor, is the quantity every drug-use claim in this thesis is anchored to: the drug-agnostic
comparison is not leak-immune, and a drug-agnostic mean baseline scores $\NIR = 0.399$ when the
comparison set may span plates against $0.161$ when it may not, precisely because it inherits that
leakage. A refinement matters here and is
developed in Section~\ref{sec:methods-controls}: if the swapped-in drug happens to have a
\emph{similar} true response then an unchanged output is correct rather than blind, so the swap
partner is chosen from three defined similarity strata and the evidence is the \emph{gradient} across
them rather than any single value.}
\label{fig:scramble}
\end{figure}


The other two instruments are \emph{forced-choice grading} --- a prediction is scored against its
own condition's truth and against another drug's truth from the same comparison set, recording how
often the own truth wins; chance is $0.50$ by construction, and the ceiling, the same forced choice
with a real held-out replicate substituted for the prediction, lands between $0.67$ and $0.83$
depending on cells per condition --- and \emph{output invariance}, which asks whether the drug token
changes the generated text at all, without reference to any truth. For a fixed control cell, the
model generates twice from the \emph{same} prompt and once from each of several prompts differing
only in the drug, and the instrument compares
\begin{equation}
  \textrm{gap} \;=\; \underbrace{\overline{\textrm{sim}}(\textrm{same prompt, resampled})}_{\textrm{sampling noise alone}}
  \;-\; \underbrace{\overline{\textrm{sim}}(\textrm{prompts differing in the drug})}_{\textrm{noise plus any drug effect}},
\end{equation}
with similarity measured as top-$100$ Kendall $\tau$ and as Jaccard overlap. A positive gap means
two generations for the same drug resemble each other more than generations for different drugs
do; a gap of zero means swapping the drug perturbs the output no more than resampling at the same
temperature does. Reported over $8$ contexts $\times$ $12$ drugs at temperature $0.8$. This
instrument is the weakest of the four --- it can only detect whether the output moves, not whether
it moves \emph{correctly} --- but it is also the one least able to be confounded, since it needs no
ground truth at all.

\begin{table}[h]
\centering
\begin{tabular}{lrl}
\toprule
Instrument & Value & Reading \\
\midrule
Forced-choice grading & $0.48$ ($\approx$ scramble; ceiling 0.67--0.83) & chance \\
Output invariance (gap) & \ci{0.000}{-0.019}{+0.016} & null \\
Scramble $\DEdr$, tier 1 & $0.740$ real vs $0.739$ scrambled & identical \\
Scramble $\DEdr$, tier 2 & $0.733$ real vs $0.729$ scrambled & identical \\
$\gap$ (identifiable subset) & \ci{+0.014}{-0.016}{+0.042} & null \\
\bottomrule
\end{tabular}
\caption{Four independent instruments agree. Near-ceiling $\DEdr$ is achieved \emph{identically}
with a wrong-mechanism drug, so the score is drug-independent.}
\label{tab:blind}
\end{table}

Before asking what the model knows, it is worth asking what there is to know. A ladder of
mean-shift baselines --- each predicting the average response of a progressively narrower group ---
answers that directly. On tier 2, predicting the global mean scores $\DEdr$ $0.056$; conditioning
on cell line raises it roughly two-and-a-half-fold to $0.142$; conditioning on mechanism does not move it ($0.052$);
and conditioning on mechanism \emph{and} cell line is no better than cell line alone ($0.120$
against $0.142$). Tiers 1 and 3 tell the same story. The informative conditioning variable in this
data is therefore the cellular context, not the compound class --- worth holding onto, because
\S\ref{sec:res-mechanism} finds precisely that asymmetry inside the model's activations, and
\S\ref{sec:res-scope} finds that mechanism does carry signal once the generic program is
subtracted. The apparent tension is a frame difference rather than a contradiction: in
full-profile space the generic response dominates and swamps a mechanism effect that becomes
visible only in the residual frame.

On the calibrated metric the model sits at chance and is matched by predictors that know nothing
about the drug: within plate on tier 2 ($n = 606$), the replicate ceiling is $0.576$, the model
scores $0.498$, a drug-agnostic linear map $0.500$, a zero-drug-information control-copy $0.504$,
and the global mean $0.180$. A separate and stronger observation, which belongs to the argument
rather than to plate leakage, is that \emph{within} plate and on the identifiable stratum a
control-copy baseline reaches $0.766$ against the model's $0.768$.

Three objections to this verdict close on measurement rather than argument. First, \emph{is the
task winnable at all?} Only $46.2\%$ of (drug, cell line) conditions are identifiable in principle;
$27.3\%$ of tested conditions are statistically inert against DMSO and $78.7\%$ have a plate-mate closer than
their own replicate. But restricting to the identifiable subset does not rescue the model: it
appears to beat the linear baseline there ($0.768$ vs $0.531$), yet a zero-drug-information
control-copy matches it exactly ($0.766$) and the leak-immune $\gap$ remains null. The model is
drug-blind \emph{even where the task is provably winnable}. The ceiling itself rises with
aggregation ($0.61$ at 10 cells, $0.76$ at 40, $0.85$ at 120), so the signal is real and
noise-limited.

\begin{figure}[htbp]
\centering
\includegraphics[width=\textwidth]{fig-difficulty}
\caption[Difficulty structure and the stratum ladder]{\textbf{The task is winnable where the noise
ceiling is high --- and precisely there, a predictor containing no drug information matches the
model.} (a) The ceiling rises steadily with the number of cells per condition, cross-plate and
within plate, so identifiability is noise-limited rather than absent: the drug signal is real and
emerges under aggregation. (b) Conditions binned by their own ceiling, with five arms in every
stratum. On the identifiable stratum ($n=204$) the model reaches $0.768$ --- and a control-copy
baseline that simply reproduces the untreated cell, and therefore knows nothing whatever about the
drug, reaches $0.766$. The paired contrast beneath each stratum is
\ci{+0.002}{-0.026}{+0.028} there, clustered over cell lines. Showing the model alone across these
strata would make the visual claim that it does better on winnable drugs; the control arm is what
shows that it does not.}
\label{fig:difficulty}
\end{figure}

Second, \emph{does the representation throw the drug away?} Cell sentences keep only the rank order
of genes and discard expression magnitude, and if the drug lived in the magnitudes the diagnosis of
the following sections would be attacking the wrong thing. The objection is testable without any
model, by asking whether \emph{true} expression discriminates drugs any better than rank does. It
does not: at single-cell resolution, cosine, Pearson and Spearman similarities on true normalised
expression all sit within $0.04$ of chance ($0.529$--$0.531$ [$0.515$, $0.543$]) --- the same
near-chance regime as rank. What recovers the signal is aggregation, not representation: at
fifteen cells the control-referenced shift cosine reaches $0.793$ [$0.775$, $0.809$]. The drug is
not hiding in the magnitudes, which matters twice over. It rules out the simplest alternative
explanation for everything that follows, and it means the repair of \S\ref{sec:res-encoding} cannot
be read as ``restoring the magnitudes'': that repair keeps rank throughout and changes only
\emph{what is being ranked}. The same measurement bounds the task from below --- at single-cell
resolution the same-drug versus different-drug gap is only $+0.006$ \nolinebreak[3]$[+0.003,
+0.008]$, so whatever a model can learn here,
it cannot be learned from a single cell.

Third, \emph{does the scramble swap in similar drugs?} The swap-distance sweep resolves this:
$\gap$ is null in every distance bin including the far tail ($-0.019$), across a $2.3\times$ spread
in swap distance, with a trend correlation of $+0.051$. Telling the model it is a drug whose true
response is maximally different still does not change the output.

\begin{claim}
On the calibrated metric the model is drug-blind, and even where the task is provably winnable: at
chance overall ($0.498$ against a $0.576$ replicate ceiling), matched exactly by a
zero-drug-information control-copy on the identifiable stratum ($0.768$ against $0.766$), and
unmoved even when told the most dissimilar drug on offer.
\end{claim}

\subsection{Is the drug represented at all?}
\label{sec:methods-probe}

A null result invites two readings, and they call for opposite repairs. Either the drug never
enters the model's computation --- an encoding failure, which target- and input-side fixes might
reach --- or it enters and is ignored --- a readout failure, which only an objective-side fix could
reach. Every behavioural measurement so far is consistent with both, because both predict the same
output. The previous section has already closed a third possibility --- that the task itself cannot
be done: the ceiling rises with aggregation, a third of conditions are identifiable in principle,
and the model is at chance even on those. So the failure is the model's, and the useful question
becomes \emph{where} in the model it lives. That question cannot be answered behaviourally; it
needs the weights, and it decomposes into three sub-questions, each with its own instrument,
because each instrument's blind spot is exactly what the next must close. Is drug identity
\emph{present} in the activations? Does the network \emph{use} the region where it is written? And
when it uses that region, does it use \emph{which drug} is written there, or only that something
is? This section answers the first; the next two take the others in turn.

The first instrument is a \emph{linear decode probe}, and its logic is deliberately weak: if even a
simple linear reader can recover the drug from the activations, then the information is present in
a form the network could in principle use --- no claim is made yet about whether it does. The
construction uses forward passes only, no generation and no training of the model itself. Real
tier-2 evaluation prompts are grouped by drug; for each, the residual-stream activation is captured
at the \emph{last prompt position} --- the position from which response generation begins --- at
every layer. A cross-validated multinomial logistic-regression classifier then learns to name which
of the twelve drugs the prompt named, and its accuracy is compared against the floor measured on
shuffled labels ($1/12 \approx 8\%$).

Drug identity reaches the representation and survives in depth: the probe reads $82\%$ at layer 9
and $76\%$ at layer 16 against the $\approx 8\%$ floor. The between-drug to within-drug variance
ratio collapses about three-fold across the final layers --- the drug is present throughout, but
geometrically de-emphasised as the output approaches. Decodability alone, however, cannot carry the
argument: a probe is an external reader with access to information the downstream layers may never
consult. Whether the network \emph{itself} reads this region is a causal question, and it is the
next section's.

\begin{claim}
Drug identity is linearly readable from the residual stream deep into the network --- $82\%$
against an $8\%$ chance floor at the peak layer --- so the behavioural null of
\S\ref{sec:res-blind} is not an encoding failure. The drug is in the model; whether anything inside
the model reads it is the next question.
\end{claim}

\subsection{Is the representation used?}

Asking whether the network uses the drug region means intervening on it, and intervening requires
localising it first. The second instrument is built in four steps, and each step exists because the
unguarded version of the same measurement misled us once already.

\emph{Step one: where the drug lives.} On $480$ tier-2 prompts --- twelve drugs, forty real cells
each --- at seven layers (2, 4, 6, 8, 9, 12, 16), the twelve per-drug mean activation vectors are
centred on the global mean and stacked, and their singular value decomposition contributes its top
ten orthonormal directions: $V_{\textrm{drug}}$, the slab of the residual stream in which drug
identity is written. A cell-line subspace is built identically and serves throughout as the
positive control, because the model demonstrably \emph{does} use cell line
(\S\ref{sec:res-blind}'s grouping ladder).

\emph{Step two: purification.} The raw drug slab is partly a batch slab. Drugs are confounded with
plate and dose by the atlas's design, and probing dose from \emph{inside} the drug subspace
recovers it at $0.90$--$0.95$ at every layer: read naively, the ``drug'' axis would partly be
measuring the batch. A context subspace is therefore built from the union of cell line, plate and
dose directions, and the drug axis is orthogonalised against all of it, giving
$V_{\textrm{pure}}$ --- the part of the drug axis not explained by batch. Every causal claim below
is made on $V_{\textrm{pure}}$, never on the raw slab.

\emph{Step three: the removal gate.} Before any ``ablating the drug changes nothing'' result can be
believed, the instrument must prove that the slab actually contains the drug. The decode probe's
accuracy is measured before and after projecting the drug subspace out, and the fraction of
above-chance signal removed must exceed $0.8$ --- otherwise the null is unfalsifiable, a
measurement of nothing. The gate passes at every layer: accuracy falls from $0.41$--$0.82$ to a
constant $0.121$ floor, removing $88$--$95\%$ of the above-chance signal. The nulls below are real
properties of the model, not broken hooks.

\emph{Step four: the causal measurement, in logit space.} The true response is teacher-forced and
its per-token log-probabilities are read off clean; a forward hook then projects the slab out of
the residual stream at \emph{every} position ($h \leftarrow h - (hV)V^{\top}$) and the
log-probabilities are read again; the measure is the mean KL divergence from clean to ablated over
the response tokens, on $60$ prompts. Scoring in logit space rather than by generating is a
deliberate repair of the project's first causal attempt, which injected a steering vector and
generated: the steering vector was small against the residual-stream norm, its effect on the output
sat under the sampling-noise floor, and its ``no effect'' was uninterpretable. The ablation is run
for $V_{\textrm{pure}}$ (headline), for a random subspace of matched dimension (the null), and for
the cell-line slab (the positive control).

\begin{table}[h]
\centering
\begin{tabular}{lrrrrr}
\toprule
Layer & Decode & Decode & Ablate pure & Random & Ratio \\
      & (raw)  & (context removed) & drug (KL) & matched dim & \\
\midrule
2  & 0.408 & 0.354 & 0.155 & 0.012 & $12.7\times$ \\
4  & 0.585 & 0.429 & 1.563 & 0.084 & $18.6\times$ \\
6  & 0.692 & 0.529 & 1.435 & 0.251 & $5.7\times$  \\
8  & 0.779 & 0.569 & 0.596 & 0.042 & $14.1\times$ \\
9  & \textbf{0.821} & 0.602 & 0.418 & 0.100 & $4.2\times$  \\
12 & 0.754 & 0.873 & 0.133 & 0.035 & $3.8\times$  \\
16 & 0.758 & \textbf{0.883} & 0.015 & 0.004 & $3.6\times$  \\
\bottomrule
\end{tabular}
\caption{All seven layers pass the removal gate (88--95\% of above-chance drug signal removed);
chance is $0.083$. Two decode series are shown because they behave differently and the distinction
matters: the \emph{raw} probe peaks at layer 9 and then declines, whereas with cell line, plate and
dose projected out decodability rises monotonically to layer 16. The identity test is reported
separately below, on a later and differently controlled design; the ablation ratios here are a
statement about \emph{functional variance} and carry no claim that drug identity is read.}
\label{tab:workspace}
\end{table}

\begin{figure}[htbp]
\centering
\includegraphics[width=\textwidth]{fig-mechanism}
\caption[The drug is encoded more with depth and never read for identity]{\textbf{The model computes
an increasingly explicit drug representation and sets it aside.} (a) With cell line, plate and dose
projected out, drug decodability rises monotonically with depth, from $0.354$ at layer 2 to $0.883$
at layer 16, against a chance rate of $0.083$; the causal variance share of that same purified
subspace stays flat near $0.03$, against roughly $0.6$ for cell line. The raw probe is drawn too and
behaves differently --- it peaks at layer 9 and declines --- because the two series are easily
conflated and the dissociation rests on the context-removed one. (b) Ablation of the purified drug
subspace against a matched-\emph{dimension} random subspace, on a log axis because late-layer
divergences are small for everything. This is a functional-variance comparison: it establishes that
the subspace is not inert, and it makes no claim about whether drug identity is read, which requires
the substitution test reported in the text.}
\label{fig:mechanism}
\end{figure}

What the ablation shows is a dissociation, and Table~\ref{tab:workspace} quantifies it. With cell
line, plate and dose projected out, drug decodability \emph{rises} with depth, from $0.354$ at
layer 2 to $0.883$ at layer 16 --- the representation becomes more explicit, not less. Yet the
purified drug subspace's variance share stays flat at $\approx 0.03$, against $\approx 0.6$ for
cell line, and ablating it costs $3.6$--$18.6\times$ the matched-dimension random null across
layers. That effect is real. It is also small beside the cell-line positive control, which the same
ablation reads at $1.0$--$11.2\times$ the drug's: removing the drug subspace costs $9$--$37\%$ of
what removing cell line costs at layers 2--9, rising to $98\%$ at layer 12 and $25\%$ at layer 16.
In variance-share terms the gap is wider still, cell line carrying $15$--$20\times$ the purified
drug slab. The instrument is not the limit; the drug is a bit player.

This is also where ablation reaches its limit as an instrument. It proves the slab carries
\emph{functional variance} that downstream layers consume; it cannot prove that the slab's
\emph{drug-specific content} is what is read --- the purified directions could carry generic
strong-program energy that merely co-varies with drug across twelve compounds. The ablation ratios
in Table~\ref{tab:workspace} therefore stand as a statement about functional variance, and what
survives from the pre-revision analyses stands with them: the ratios (whose null, a
matched-\emph{dimension} random subspace, was always correctly constructed), the variance shares,
and the decodability-versus-depth dissociation. Whether the network reads \emph{which drug} the
slab encodes requires replacing what the slab says rather than removing it --- and that
substitution test is the next section's.

\begin{claim}
The drug representation is used, but at a bit player's weight: decodability rises monotonically
with depth while the slab's causal share stays flat at $\approx 0.03$, roughly $5\%$ of the cell
line's $\approx 0.6$. Ablation cannot say whether the slab's content is read --- that needs a substitution.
\end{claim}

\subsection{Is the drug's identity read?}
\label{sec:res-mechanism}

The substitution test is the only instrument that can answer this, and getting one to survive audit
took three failed designs. They are stated rather than hidden, because each failure teaches a
constraint the final design embodies --- and because the first two produced results this thesis
previously quoted as findings.

\emph{Failed design one: mismatched geometry.} The first test overwrote drug A's code with drug
B's, inside a $\leq 10$-dimensional purified subspace, and compared the resulting KL against a
random vector of identical norm drawn in the \emph{full} $2048$-dimensional stream --- so roughly
$99.5\%$ of the control's energy lay in directions the intervention never touched, including the
cell-line directions this chapter measures at $\approx 0.6$ variance share. ``The swap moves the
output less than noise'' is exactly what that geometry predicts, whatever the readout does.

\emph{Failed design two: ablate-then-add.} The second drew the control inside the subspace, which
looked like a repair and was not: the hook ablates before adding, so what is delivered is (added
$-$ removed), and an uncentred class mean shares its global-mean component with the thing removed.
The substitution arm was then close to a restoration and the control arm a real displacement.
Simulated with no model present, the control delivers $5.3\times$ the perturbation in $100\%$ of
draws, with the predicted gap $2\langle c_B, s\rangle = 9.179$ against an observed $9.162$. The
retracted result was geometry, not biology.

\emph{Failed design three: the normaliser.} The third injected the displacement and scored
$\Delta \log P(r_B)$ against a norm-matched random direction --- and was confounded the opposite
way. A random direction inside the slab is off the class-mean manifold: it inflates the logsumexp
normaliser (measured $+2.75$ nats) where the on-manifold displacement does not, so \emph{every}
target is depressed under the control arm whether or not it carries identity; with the scored
target stripped of all identity information, $36\%$ of the effect still stood. Norm-matched is not
disruption-matched. The same argument retires the simplest test of all, ``how far did the output
move'': distance moved is exactly what norm-matching equalises, so that test has no power in
precisely the world where the readout does read the drug.

The design that survived makes two structural changes so that the arms differ in nothing but which
drug they name, and it ablates nothing. First, a \emph{contrast} score,
\begin{equation}
  \log P\!\left(r_B \mid \textrm{prompt}_A\right) \;-\; \log P\!\left(r_A \mid \textrm{prompt}_A\right),
\end{equation}
where $r_B$ is a real response emitted under drug B by a held-out cell. Any shift common to both
responses --- the entire normaliser included --- cancels algebraically. What survives is the
model's \emph{preference} between drug B's own held-out response and drug A's. Second, every arm is
anchored at drug A as an on-manifold class-mean displacement of identical norm: the \texttt{swap}
arm injects $P(\mu_B - \mu_A)$, pointing at the drug whose response is scored; the \texttt{permuted}
arm injects $P(\mu_C - \mu_A)$, rescaled to identical norm --- the headline control, same anchor,
same construction, wrong drug; and an isotropic random direction in the slab is kept only as an
off-manifold damage diagnostic, never as the comparison, because a random direction disrupts the
model in ways an on-manifold class-mean displacement does not. The subtraction $\mu_B - \mu_A$
cancels the global mean exactly, so the injected vector is pure drug difference, and no ablation
anywhere means there is no ``removed'' term for the control to trip over.

Everything is out of sample, and everything is on a comparable scale. The subspaces and class means
are fit on one slice of cells; $r_B$ is drawn from a third disjoint slice --- never from the rows
defining $\mu_B$, else the injected vector and the scored response would share cells and the test
would be partly circular; all scoring happens on the held-out slice. And because the five training
arms emit different target formats, raw log-probability differences live on different scales, so
every effect is divided by that model's own \emph{clean} preference gap between A's and B's
responses, measured with no intervention: the reported unit is the fraction of the existing
preference gap the injection closes, comparable across arms.

The guards around the measurement are the point of the exercise, so they are listed in full.
Uncertainty is a cluster bootstrap over ordered (A,B) pairs --- rows sharing a pair share both the
injected vector and the scored response, which is the dependence the resampling must respect.
Equivalence is by TOST: the interval must lie inside a pre-registered margin \emph{and} contain
zero, with the margin set at $10\%$ of the model's own clean preference gap. A per-layer headroom
gate excludes any layer where the cell-line positive control shows less than $5\times$ dynamic
range, because a saturated instrument can support neither verdict. A drug-proxy screen refuses any
context label that is a near-bijection with drug --- it caught the metadata field
\texttt{dose\_float} holding a sample identifier in the residual builders, which, had it been
orthogonalised against, would have deleted the drug axis by construction. Multiplicity is handled
by Benjamini--Hochberg across depths. A magnitude ladder repeats the test at $\alpha \in \{0.5, 1,
2, 5, 10\}$ times the natural displacement, because a real identity effect should show a monotone
dose--response. An injected-norm diagnostic reports the displacement against the cell's
\emph{centred} slab component ($0.69$--$1.40$: the injection is the size of the thing it displaces,
so a null cannot be blamed on a whisper). And a selftest with teeth runs three planted worlds ---
drug inert, drug read, and drug inert behind a norm-sensitive readout --- through the actual
headline code path; it separates them ($-0.042$, $+1.520$, $-0.010$), and it failed four times
during development, each failure a real defect that the list above now guards.

This design was run across five training arms and, for the two residual arms, replicated at three
further seeds. Of $26$ headline tests --- five arms by up to seven depths --- exactly one survives
Benjamini--Hochberg correction: the residual arm at layer $12$ ($q = 0.026$), with a monotone
dose--response across the ladder's first three rungs ($+0.0013$, $+0.0017$, $+0.0044$ at
$\alpha = 0.5, 1, 2$). The consensus arm at layer 2 ($p = 0.012$) does not survive ($q = 0.156$).

\begin{table}[h]
\centering
\begin{tabular}{lrrr}
\toprule
Arm / seed & effect & 95\% CI & $p$ \\
\midrule
residual, seed 42 & $+0.0197$ & $[+0.0063, +0.0335]$ & $0.0010$ \\
residual, seed 43 & $+0.0118$ & $[+0.0060, +0.0173]$ & $0.0003$ \\
residual, seed 45 & $+0.0085$ & $[+0.0029, +0.0143]$ & $0.0015$ \\
\midrule
single-cell arm, six depths & \multicolumn{3}{l}{nothing detectable at any depth (median $-0.0012$)} \\
\bottomrule
\end{tabular}
\caption{The identity test, as a fraction of that model's own clean preference gap between the two
responses being compared (mean effect $+0.0133$ across the three seeds that measured it). Seed 44
produced no layer-12 measurement --- it lost five of seven layers to the headroom gate --- so the
pre-registered criterion of two replications in three new seeds is met with the third a
non-measurement rather than a failure.}
\label{tab:identity}
\end{table}

The value of this is that it is a \emph{positive control on real data}. The residual arm is
independently known to use the drug: it is the arm whose prompt sensitivity is established
behaviourally in \S\ref{sec:res-generalisation}. The probe detects identity reading there,
replicably, and detects none anywhere in the arm that is behaviourally drug-blind. A null from an
instrument that has never been shown to register a positive is uninterpretable; this one has.

\begin{claim}
The generation readout consults the drug representation \emph{at very low gain}: about $1\%$ of the
preference gap, an order of magnitude below the pre-registered margin for calling an effect
material, in the one arm where it is detectable at all. The claim is not that the readout cannot
read drug identity, and it is not the stronger mechanistic story that earlier drafts of this thesis
told.
\end{claim}

Three limits belong with it. The effect appears at one depth of seven, and layer 4 is positive in a
single seed and null in the other three, which is what multiplicity discipline is for. The
single-cell arm has no layer-12 measurement at all --- its cell-line positive control reached
$3.88\times$ there, under the $5\times$ headroom gate, so the layer was excluded --- meaning the
cleanest comparison, the same depth across two arms, does not exist; the dissociation rests on that
arm being null at all six depths it could measure. And the sibling residual-holdout arm does not
reproduce the effect ($+0.0028$ and $+0.0049$ in the two seeds where it was measured). Read
together with the decodability series, the picture is a representation that becomes more explicit
with depth while the process that writes the output barely consults it --- which is suggestive of
the workspace account of Section~\ref{sec:lit-interp}, and is offered as an interpretation of the
pattern rather than as a demonstration of the mechanism.

\subsection{No fix that keeps the encoding works}
\label{sec:res-epsilon}

The mechanism makes a falsifiable prediction, and this section is its test. If the defect is a
readout that feels the magnitude of activity in the drug region but not its direction, then
improving what sits \emph{in} that region cannot help --- and improving the training target is
exactly that kind of intervention. The obvious objection is that one failed target construction
proves nothing, so rather than try another we swept the family. Two constructions anchor the
sweep. \emph{Consensus} groups cells by (drug, cell line, dose) and averages them into one
pseudobulk sentence, with every original prompt retained so only the response changes.
\emph{Optimal transport} computes, for each (cell line, plate) group, a coupling $\pi$ between
control and treated cells minimising $\sum_{ij}\pi_{ij}\lVert x_i - y_j\rVert^2$ under uniform
marginals, solved by Sinkhorn iteration~\cite{cuturi} on $K = \exp(-C/\varepsilon)$; couplings are
computed in the released Tahoe scVI latent~\cite{scvi} ($n_{\textrm{latent}} = 10$), joined per
cell by barcode rather than re-encoded --- a re-encode produced a degraded latent (cell-line probe
$0.11$ versus $0.975$ for the shipped latent) --- and targets are built \emph{decoder-free}, as
convex combinations of real treated cells in expression space, so they always lie on the data
manifold. The regularisation parameter then defines a ladder whose endpoints are experiments we had
already run: $\varepsilon\rightarrow\infty$ gives a uniform coupling and hence the treated mean
(consensus); $\varepsilon\rightarrow 0$ gives near-hard assignment, approximately a single
state-matched treated cell. This is what makes the sweep conclusive rather than anecdotal.

\begin{table}[h]
\centering
\begin{tabular}{llrl}
\toprule
Target & $\varepsilon$ & $\gap$ & Verdict \\
\midrule
single cell (real treated cell) & $\approx 0$ & \ci{+0.014}{-0.016}{+0.042} & null \\
OT barycentric (state-matched) & $0.5$ & \ci{-0.001}{-0.018}{+0.016} & null \\
consensus (global mean) & $\rightarrow\infty$ & \ci{-0.010}{-0.022}{+0.003} & null \\
\bottomrule
\end{tabular}
\caption{The regularisation parameter interpolates between the two target constructions tested
independently, and no point on the ladder lifts the model out of the near-null regime. These figures
come from a \emph{random-partner} scramble; \S\ref{sec:res-encoding} shows that under a sharper
stratified test the optimal-transport arm does show a weak effect
(\ci{+0.0263}{+0.0069}{+0.0476}), roughly five-fold below what re-encoding the target achieves.
The consensus endpoint is additionally the ladder's weakest-evidenced point: that arm was stopped at
roughly $60\%$ of an epoch and its tier-1 NIR was unscoreable, so the negative rests on the
single-cell and optimal-transport endpoints.}
\label{tab:epsilon}
\end{table}

Training was not inert: the consensus model stops echoing the control ($0.653$ against a
control-copy of $0.766$, where the single-cell model sits \emph{on} the control at $0.771$) and its
predictions are not mode-collapsed. But the variation is uncorrelated with real drug identity
(Mantel $\approx +0.03$). Not collapse --- wrong structure.

\begin{claim}
Denoising changes \emph{what is in} the target while the defect is what its tokens encode; that is
the quantity the next two sections define and then measure.
\end{claim}

\subsection{The residual frame: predicting only what makes this drug different}
\label{sec:methods-residual}

If the encoding is the constraint, the target must be redefined so that its tokens \emph{are} the
drug-specific part. For each condition (drug, cell line, plate, dose) with at least $40$ treated
and $20$ plate-matched control cells:
\begin{equation}
  \resid(d,c) \;=\; \underbrace{\big(\bar{y}_{\textrm{treated}} - \bar{y}_{\textrm{control}}\big)}_{\shift(d,c)}
  \;-\; \underbrace{\frac{1}{|D_c|}\sum_{d' \in D_c} \shift(d',c)}_{\textrm{generic}(c)} .
\end{equation}
The control subtraction removes the cell's own state; the mean-over-drugs subtraction removes the
\emph{generic drug program} --- the stress and cell-cycle response common to all compounds, which
\S\ref{sec:res-metrics} showed accounts for essentially all of the model's apparent skill. What
remains is what makes \emph{this} drug different. Four properties of how the generic is estimated
decide whether that remainder can be trusted, and each was forced by a measured failure.

First, \textbf{the split precedes the fit}. An earlier build estimated the generic and the
reliability filter over \emph{every} condition and assigned the holdout afterwards, so each held-out
target was defined partly by other held-out conditions. The pipeline is therefore staged ---
inventory (metadata only), split (metadata only, before any pseudobulk exists), fit (the generic is
handed the training keys explicitly), transform (held-out conditions are projected through the
fitted generic and never enter it) --- and a poison test stands behind the staging: corrupting every
held-out expression vector must leave the training targets byte-identical, and corrupting a training
condition must change them. What the split is, and why it is keyed at whole wells, is the subject
of \S\ref{sec:methods-holdout}; here it matters only that it exists and is metadata-only.

Second, \textbf{the scope is the plate, not the cell line}. The original argument for cell-line
scope was reproducibility --- $62\%$ of conditions reproduce across a half-split at cell-line scope
against $19\%$ at plate scope, on the reasoning that a plate holds too few drugs to estimate a mean
and subtracting a noisy one injects noise. That comparison does not hold up: both figures were
computed with a \emph{shared} control pseudobulk subtracted from both halves, which makes the halves
correlate through their common control error, and the inflation is not neutral between scopes ---
it favours cell-line scope precisely because the cell-line generic is itself less noisy. Measured
with disjoint control halves, the two are comparable: $20\%$ at plate scope against $16\%$ at
cell-line scope. Two further diagnostics settle it in plate's favour and are reported in
\S\ref{sec:res-transfer}: at cell-line scope the dose-transfer ordering inverts and the
shared-minus-split control bias reads $+0.254$, against $-0.000$ at plate scope. A plate group with
fewer than three other training drugs yields no generic, and the condition is dropped and counted
rather than silently falling back to the contaminated frame.

Third, \textbf{the mean is over drugs, not conditions}, and \textbf{leaves the condition's own drug
out}. A drug measured at five doses would otherwise contribute five times the weight of a
single-dose drug. And once the generic is train-only, an asymmetry appears that had not existed
before: a training condition's generic would contain its own drug while a held-out condition's would
not, so the two splits would be scored against differently defined targets and the generalisation
gap would partly measure that difference. The generic for every condition is therefore ``the mean
over training drugs other than this one, within scope'' --- identical in definition on both sides.

Fourth, \textbf{the reliability line is resolved from the data, not inherited}. A condition is kept
for training only if its residual reproduces across a half-split, $\cos(\resid_A, \resid_B)$ above
a threshold --- with each half measured against its \emph{own} control half, so the two share
nothing. The threshold itself was the single largest determinant of the training set, and it had
been inherited ($0.2$) from an era when reliability was measured with shared controls --- a
different quantity on a different scale, under which the mean is $\approx 0.24$ rather than
$+0.127$. Rather than pick a new number, the builder reports the whole retention curve against a
null that pairs half A of one condition with half B of a \emph{different} condition --- same
encoding, same dimensionality, same noise, no shared biology --- and the threshold is set at the
null's $95$th percentile: a condition is kept when its two halves agree more than two unrelated
conditions' halves do. At the measured null ($95$th percentile $+0.109$) this retains $2{,}523$ of
$4{,}999$ training conditions ($50.5\%$). The held-out sets are \emph{not} filtered by this
criterion: dropping held-out conditions for failing their own reproducibility bar would select the
evaluation set on its own outcome, so noise in retained training conditions costs learning
efficiency, not validity.

What this fraction means is stated carefully, because it is easy to over-read.

\begin{claim}
The half-split agreement is a \emph{within-well sampling-precision} criterion: it asks whether two
disjoint cell subsets of the \emph{same} well give the same residual. It is not biological
reproducibility --- $286$ of $287$ distinct (drug, dose) treatments in this atlas occur in exactly
one well, so no independent-well replicate exists to measure that with --- and the $50.5\%$ is a
statement about how much drug-specific signal survives pseudobulk noise at this cell budget, not
about how much biology is real.
\end{claim}

The target string is \texttt{<up genes> \DOWN{} <down genes> \ENDCELL}, $100$ genes per block,
ordered by signed residual magnitude; sign is biology, so it is given a symbol rather than being
folded into a magnitude ranking. Predicted and true residuals are both mapped to a signed rank
vector (top-$k$ up positive, top-$k$ down negative, weight $1/\log_2(\textrm{rank}+1)$) so a
generated sentence and a continuous truth are compared like for like; similarity is cosine and the
score is $\NIR$ over other drugs in the same cell line.

% fig-frame.tex -- the residual frame and the encoding fork.
% No data dependency. Annotations from RESULTS/target_divergence.json, cellline scope:
%   full 34.59 / shift 111.42 / residual 118.25 ; ceilings 0.7420 / 0.7420 / 0.7432.
% Layout notes:
%   * two subtraction ROWS, not an L-shape (the L put the generic term on a band that collided
%     with its own caption).
%   * the caption for the generic sits BELOW both rows, not between them -- between them it ran
%     under the divider and into the right-hand column.
%   * row labels are centred on profiles spaced 1.65cm apart; at 1.45cm "generic(c)" and
%     "residual r(d,c)" touch.
\begin{figure}[htbp]
\centering
\begin{tikzpicture}[
  x=1cm, y=1cm, font=\small,
  bar/.style={fill=black!55, draw=none},
  faintbar/.style={fill=black!22, draw=none},
  gene/.style={draw=black!35, rounded corners=1pt, inner sep=1.8pt, font=\ttfamily\scriptsize},
  shared/.style={gene, fill=black!5, text=black!45},
  differs/.style={gene, fill=black!16, text=black, draw=black!70, thick},
  lbl/.style={font=\scriptsize, text=black!70},
  strong/.style={font=\scriptsize, text=black!90},
  op/.style={font=\normalsize, text=black!75},
]
\newcommand{\prof}[4]{%
  \foreach \h [count=\i from 0] in {#3} {
    \fill[#4] ($(#1,#2)+(\i*3.0pt,0)$) rectangle ++(2.3pt,\h pt);
  }%
}

% ================================================ LEFT: two subtractions
\node[lbl]    at (0.52,3.34) {treated};
\node[lbl]    at (2.17,3.34) {control};
\node[strong] at (3.90,3.34) {shift $\shift(d,c)$};
\prof{0.10}{2.62}{17,12,9,14,7,10,5,8}{bar}
\node[op] at (1.48,2.92) {$-$};
\prof{1.75}{2.62}{10,8,6,11,4,7,3,6}{bar}
\node[op] at (3.13,2.92) {$=$};
\prof{3.40}{2.62}{7,4,3,3,3,3,2,2}{bar}

\node[strong] at (0.52,1.94) {shift};
\node[lbl]    at (2.17,1.94) {generic$(c)$};
\node[strong] at (3.95,1.94) {residual $\resid(d,c)$};
\prof{0.10}{1.22}{7,4,3,3,3,3,2,2}{bar}
\node[op] at (1.48,1.52) {$-$};
\prof{1.75}{1.22}{6,4,3,2,3,2,2,2}{faintbar}
\node[op] at (3.13,1.52) {$=$};
\prof{3.40}{1.22}{1,0,0,1,0,1,0,0}{bar}

\node[lbl, anchor=north, align=center] at (2.30,0.82)
  {generic$(c)$ is the mean over the \emph{other}\\drugs in this cell line};

\draw[black!25] (5.05,0.10) -- (5.05,3.60);

% ================================================ RIGHT: the encoding fork
\node[strong, anchor=west] at (5.35,3.44) {the \emph{same} measurement, written two ways};

\node[lbl, anchor=west] at (5.35,2.96) {\textbf{(a)} genes ranked by \emph{expression}};
\node[anchor=west] at (5.35,2.52) {%
  \tikz[baseline]{\node[shared]{ACTB};}\;\tikz[baseline]{\node[shared]{GAPDH};}\;%
  \tikz[baseline]{\node[differs]{CDKN1A};}\;\tikz[baseline]{\node[shared]{TPT1};}\;$\cdots$};
\node[anchor=west] at (5.35,2.06) {%
  \tikz[baseline]{\node[shared]{ACTB};}\;\tikz[baseline]{\node[shared]{GAPDH};}\;%
  \tikz[baseline]{\node[differs]{HMOX1};}\;\tikz[baseline]{\node[shared]{TPT1};}\;$\cdots$};
\node[strong, anchor=west] at (5.35,1.58)
  {only \textbf{34.6 / 200} tokens differ \;(\textbf{17\%})};

\node[lbl, anchor=west] at (5.35,1.00) {\textbf{(b)} genes ranked by \emph{signed residual}};
\node[anchor=west] at (5.35,0.56) {%
  \tikz[baseline]{\node[differs]{CDKN1A};}\;\tikz[baseline]{\node[differs]{MDM2};}\;%
  \tikz[baseline]{\node[gene, fill=black!3]{[DOWN]};}\;\tikz[baseline]{\node[differs]{CCNB1};}\;%
  $\cdots$};
\node[strong, anchor=west] at (5.35,0.08)
  {\textbf{118.2 / 200} differ \;(\textbf{59\%})};

% ================================================ the invariant
\draw[decorate, decoration={brace, amplitude=4pt}, black!55] (0.10,-0.55) -- (11.20,-0.55);
\node[strong, anchor=north] at (5.65,-0.73)
  {and the replicate-retrieval ceiling is \textbf{unchanged}: $0.742 \rightarrow 0.743$};

\end{tikzpicture}
\caption[The residual frame and the encoding fork]{\textbf{The information was always present; the
tokenisation discarded it.} Left: the drug-specific residual is built by two subtractions. The
control removes the cell's own state, leaving the \emph{shift}; the mean over the other drugs in the
same cell line removes the \emph{generic} program that every compound triggers, leaving what makes
this drug different. Right: that same measurement, written two ways. Ranking genes by expression (a)
puts housekeeping genes first, so only $34.6$ of the top $200$ tokens differ between two drugs in the
same context and cross-entropy has almost nothing to learn drug identity from; ranking by signed
residual (b) raises that to $118.2$. The bracket is the point of the figure --- the
replicate-retrieval ceiling is essentially identical under both encodings, so no information has been
added. What changed is how much of it the token sequence exposes to the loss.}
\label{fig:frame}
\end{figure}


\subsection{How much signal is there to encode}
\label{sec:methods-vardecomp}
\label{sec:res-transfer}

Before training anything on the residual, it is worth knowing how much drug-specific signal exists
and what kind it is, because the answer scopes every claim in this chapter. The shift's energy
decomposes exactly into residual $0.621$, generic $0.390$ and a cross term of $-0.011$: $62.1\%$ of
the perturbation response's energy is drug-specific, $39.0\%$ is the generic program, and the two
components are essentially orthogonal, so the shares are genuine. Every claim that follows is
scoped by the first of those numbers, and it has not previously been computed for this task.

Within the drug-specific part, two components call for different responses. Modelling
\begin{equation}
  \resid(d,c) \;=\; \beta(d) \;+\; \kappa(d,c) \;+\; \epsilon,
  \qquad
  \Tcoef \;=\; \frac{\operatorname{var}\beta}{\operatorname{var}\beta + \operatorname{var}\kappa},
\end{equation}
the per-drug constant $\beta(d)$ is retrievable without any model --- it is what a lookup table
reproduces --- while the drug$\times$cell-line interaction $\kappa(d,c)$ is the only component a
conditional model could win on. This is DrEval's naive predictor, in transcriptomic space: their
\texttt{NaiveMeanEffectsPredictor}~\cite{dreval} predicts $\hat{y}_{ij} = \mu^c_i + \mu^d_j - \mu$,
an additive model of cell-line and drug main effects with no interaction term, and our
$\textrm{generic}(c) + \beta(d)$ is precisely that quantity lifted from a scalar potency to a
$946$-dimensional profile. It follows that $\kappa$ \emph{is} the term their naive model omits, and
that $1 - \Tcoef$ measures exactly the headroom available above it. DrEval finds that deep models
barely beat this additive baseline on drug sensitivity; the transfer coefficient asks whether, in
the transcriptomic setting, there is anything above it to win in the first place.

$\Tcoef$ is estimated by disattenuated cross-line correlation,
\begin{equation}
  \hat{\Tcoef} \;=\; \frac{\cos\!\big(\resid(d,c_1),\, \resid(d,c_2)\big)}
                           {\sqrt{\rho(d,c_1)\,\rho(d,c_2)}},
  \qquad
  \rho \;=\; \frac{2\cos(\resid_A, \resid_B)}{1 + \cos(\resid_A, \resid_B)},
\end{equation}
with $\rho$ the Spearman--Brown reliability of the full-sample residual. Dividing out the
reliability is essential --- measurement noise alone depresses the raw cross-line cosine and would
otherwise be read as interaction --- and the reliability entering it is a \emph{within-well}
split-half precision, which is optimistic relative to true biological reproducibility. That
optimism sits in the denominator, so the interaction share $1 - \Tcoef$ is, if anything,
understated.

The measurement itself carries a control-design lesson, stated because it changed the protocol.
The first run compared $\Tcoef$ against a negative control of \emph{different drugs on the same
plate}, which returned $+0.485$ against $\Tcoef = 0.511$, and was correctly declared void: that
control differs from the estimand in two ways at once --- different drug \emph{and} same cell line
--- and plate structure surviving a cell-line-scoped generic inflates same-plate comparisons (to
$+0.478$) while leaving cross-line comparisons untouched. A \emph{structure-matched} control ---
different drug, different cell line, breaking only the drug match --- is the one that bounds the
estimand, and a negative control that differs in more than one way cannot do so. Five controls in
all accompany the estimate: the control half-split check (a shared control pseudobulk inflates
reliability and biases $\Tcoef$ downward --- measured at $+0.254$ at cell-line scope against
$-0.000$ at plate scope); the structure-matched negative control; a simulated null calibrated to
the measured noise level (at $\kappa = 0$ the estimator does not return $1.0$ --- the raw
cross-line cosine there is $0.734$, which read naively reports $27\%$ interaction where the truth
is zero); dose strata, reported separately as a reference scale with dose resolved to a
\emph{molar} quantity (an earlier cache carried the well identifier in the dose field, so what was
labelled dose transfer compared physical samples rather than concentrations); and the leave-one-out
generic of \S\ref{sec:methods-residual}. Validation against planted ground truth recovers
$\Tcoef = 1.00$, $0.70$ and $0.40$ to within $0.005$, and the different-drug negative control
returns $-0.001$.

\begin{table}[h]
\centering
\begin{tabular}{lrrr}
\toprule
Quantity & Estimate & 95\% interval & Reference \\
\midrule
$\Tcoef$ (cross cell lines, plate-scoped generic) & $0.553$ & $[0.514, 0.593]$ & simulated $\kappa{=}0$ null $\approx 1.0$ \\
drug $\times$ cell-line interaction share & $44.7\%$ & $[40.7\%, 48.6\%]$ & --- \\
$\Tcoef$ across dose (same drug, same line) & $0.707$ & $[0.632, 0.782]$ & reference scale \\
matched negative controls & $-0.017$, $-0.008$ & spanning zero & must be $\approx 0$ \\
\bottomrule
\end{tabular}
\caption{The transfer coefficient, estimated on real molar doses with dyadic cluster-robust
intervals. The matched negative control sits on zero; the same-plate control that produced the
earlier false alarm reads $+0.478$ for the reasons given in the text.}
\label{tab:transfer}
\end{table}

Roughly $44\%$ of the drug-specific residual variance is drug$\times$cell-line interaction, robust
to the generic's scope and to control splitting. Two internal checks agree, and both were decisive
between scopes. \emph{Dose ordering}: same-drug, same-line, different-dose transfer is $0.707$
[$0.632$, $0.782$], above the cross-line $0.553$ --- changing the dose costs \emph{less} than
changing the cell line, the biologically correct ordering; at cell-line scope the ordering
inverts. \emph{Control-split bias}: at plate scope the shared- and split-control builds agree to
$-0.000$, against $+0.254$ at cell-line scope, so that discrepancy was plate contamination rather
than the split. The per-drug constant $\beta(d)$ --- which the lookup of \S\ref{sec:res-lookup}
reproduces almost exactly --- is therefore \emph{not} the whole signal. There is a conditional
target worth about half the drug-specific variance that a per-drug constant cannot express.

Whether anything could \emph{learn} that interaction is a separate question. If a cell line
modulated drug responses in a consistent direction, that direction would be estimable from the
line's own cells and a conditional model would have a concrete target. Writing
$\kappa(d,c) = \resid(d,c) - \hat{\beta}(d)$, the test compares $\cos(\kappa(d,c), \kappa(d',c))$
for different drugs in the \emph{same} cell line against the same quantity across \emph{different}
cell lines, with $\hat{\beta}$ estimated \textbf{leave-one-cell-line-out} --- estimated over all
lines it would contain $\resid(d,c)$ itself, and subtracting a mean containing its own term induces
a $-1/(m-1)$ correlation that biases the signal arm downward by construction, which would make a
structured interaction look idiosyncratic, the exact error the test exists to avoid. Both arms are
disattenuated by $\kappa$'s own split-half reliability, since $\kappa$ is a difference of two noisy
quantities and is therefore noisier than the residual it comes from.

\begin{table}[h]
\centering
\begin{tabular}{lrrr}
\toprule
Configuration & same cell line & different cell lines & excess within-line \\
\midrule
plate-scoped generic & \ci{-0.004}{-0.011}{+0.002} & $+0.003$ & $-0.007$ \\
cell-line-scoped & \ci{-0.003}{-0.021}{+0.015} & $+0.002$ & $-0.005$ \\
\bottomrule
\end{tabular}
\caption{No cell line bends different drugs in a consistent direction. $4{,}000$ pairs per arm; the
signal interval at plate scope excludes anything above $+0.002$, so this is a real null rather than
an underpowered one. The cell-line-scoped row is wider ($[-0.021, +0.015]$) and carries the claim
less strongly.}
\label{tab:kappa}
\end{table}

We looked for structure in the interaction and did not find it: no cell line modulates different
drugs in a consistent direction, and a direct test of whether mechanism-matched drugs share an
interaction with a given cell line returned nothing detectable (permutation $p \approx
0.15$--$0.31$). The temptation is to read two negative results as \emph{irreducibility}, and they
do not support it. The second test bounds its own power severely: $\kappa$'s split-half reliability
in this atlas is $0.195$, so the quantity being correlated is barely estimable, and same-mechanism
compounds are so heavily co-plated that only $3$--$7\%$ of matched pairs survive a different-plate
requirement, which leaves the plate confound in place. The defensible reading is that $\kappa$'s
structure is \emph{unresolved in this dataset}. Resolving it would take more cells per condition,
or a plating design that crosses mechanism with plate --- neither of which is an analysis choice.

One measurement remains, and it is the bridge back to the model: how much of the signal reaches
the tokens. Because a cell sentence ranks genes by expression, the target strings for two drugs in
the same context are nearly identical, and the loss weights every token position equally. The
measurement counts, for pairs of drugs sharing a context, how many of the $200$ target positions
carry a different gene in the two drugs' target strings, averaged over pairs; the ceiling column
scores within-condition split-half drug retrieval through the identical representation, so the two
columns are measured in the same units.

\begin{table}[h]
\centering
\begin{tabular}{lrr}
\toprule
Representation & Inter-drug top-200 tokens differing & Replicate retrieval ceiling \\
\midrule
full profile (the standard target) & \textbf{34.6 / 200 (17\%)} & 0.742 \\
shift (treated $-$ control) & 111.4 (56\%) & 0.742 \\
\textbf{residual (drug-specific)} & \textbf{118.2 / 200 (59\%)} & 0.7432 \\
\bottomrule
\end{tabular}
\caption{$83\%$ of the standard target's tokens are identical between any two drugs in the same
context. The residual row is quoted at \textbf{cell-line} scope; at plate scope it reads $120.3$
and $0.7467$. Note the retrieval ceiling is \emph{unchanged} across all three rows: the information
was always present, and the tokenisation discarded it.}
\label{tab:divergence}
\end{table}

The natural inference is that only a small fraction of the gradient can be carrying drug identity.
That inference is a \emph{hypothesis} about the optimisation: what is measured here is the overlap
of the target gene sets, not a token-level loss or gradient attribution, and it should be read at
that strength. Its warrant is predictive --- it explains why every change to target \emph{content}
failed identically, and it predicts that changing what the tokens \emph{encode} would not.

\begin{claim}
There is a conditional target worth about half the drug-specific variance that a per-drug constant
cannot express ($44.7\%$ interaction share); no learnable structure in it is found; and re-encoding
the target as the residual exposes $59\%$ of its tokens to the drug at an unchanged retrieval
ceiling --- the information was always present, and the tokenisation discarded it.
\end{claim}

\subsection{Re-encoding produces prompt sensitivity}
\label{sec:res-encoding}

Re-encoding the target as the drug-specific residual raises the differing fraction to $59\%$ at an
unchanged retrieval ceiling, and produces the first non-null result in the project. On the earlier
target build, the stratified scramble of \S\ref{sec:methods-controls} gave a gap that grew with
swap dissimilarity --- \ci{+0.0351}{+0.0061}{+0.0656}, \ci{+0.0716}{+0.0371}{+0.1059} and
\ci{+0.1429}{+0.1112}{+0.1787} across the \texttt{near}, \texttt{orth} and \texttt{opposite}
strata --- while the single-cell and optimal-transport controls stayed flat and null at every
stratum ($-0.001$ to $+0.016$). The monotone gradient is itself the evidence: a metric artefact
would not order itself by swap dissimilarity.

\begin{figure}[htbp]
\centering
\includegraphics[width=\textwidth]{fig-repair}
\caption[The gap grows with swap dissimilarity]{\textbf{On the earlier target build, the residual
model's advantage over its own scramble grew with how dissimilar the swapped-in drug's response was;
a control scored through the identical pipeline stayed flat and null.} The evidence is the gradient,
not any single point --- a metric artefact would not order itself by swap dissimilarity. The
horizontal axis is the measured mean residual cosine between the real and swapped-in drug, which is
byte-identical across all four runs, so the positions are exact. The three solid series are at a
matched $600$-token generation budget; the residual model's $1400$-token run is drawn separately as
the dotted overlay, since the budget alone doubles the effect. Intervals are $95\%$ bootstrap
clustered over cell lines, $n=250$ conditions per point. The \texttt{near} stratum is the weakest
test by construction: where the swapped-in drug's true response is similar, an unchanged output is
correct rather than blind. These magnitudes are superseded by the rebuilt-target run of
\S\ref{sec:res-generalisation}; the figure stands as the first demonstration of the effect.}
\label{fig:repair}
\end{figure}

Three further diagnostics on that build agreed, and only the residual model passed all three: the
correlation between condition reproducibility and model $\NIR$ was $+0.144$ (the model does better
where the drug effect is real) against $-0.117$ and $-0.216$ for the controls; prediction diversity
was highest for the residual model; and while $\cos(\textrm{prediction}, \textrm{own truth})$
exceeded $\cos(\textrm{prediction}, \textrm{other drugs})$ for all three arms, the residual model's
margin ($+0.084$) was roughly twice the optimal-transport arm's ($+0.049$) and three times the
single-cell arm's ($+0.026$).

Two measurement lessons belong with this result. The effect was first measured at a $600$-token
generation cap, which silently truncated $26\%$ of generations --- a signature is $\approx 201$ gene
symbols at 2--4 BPE tokens each --- and \emph{halved} it ($+0.0716$ against the correct $+0.1429$);
re-running both controls at the matched $1400$-token budget left them flat and null, so the
comparison is budget-matched and the conclusion unaffected. And the optimal-transport arm,
originally reported as achieving nothing, scores \ci{+0.0263}{+0.0069}{+0.0476} under the sharper
stratified test with a monotone gradient: it does use the drug, weakly. The original null was an
underpowered instrument, not an absence of signal; the $\varepsilon$-ladder conclusion stands, but
``OT achieved nothing'' is withdrawn.

What those magnitudes can and cannot claim is re-adjudicated in the next section, because they were
read against the \texttt{opposite} comparator, and that comparator turns out not to be a null. The
comparator-free evidence is already in hand, and it is the strongest in this chapter: regressing
each generation's $\NIR$ on the cosine between the named drug's truth and the target's truth gives
a slope of \ci{+0.390}{+0.225}{+0.542} on trained conditions --- zero slope is the natural null, so
no comparator defect can reach it. Re-encoding the target made the model respond to the combined
drug-and-mechanism prompt. What that response is worth is the question the rest of the chapter
answers.

\begin{claim}
Re-encoding the target produced the first non-null drug effect in the project: the gap grows with
swap dissimilarity while controls stay flat, and the comparator-free slope of retrieval against the
named drug's truth is \ci{+0.390}{+0.225}{+0.542}. Sensitivity is established; fidelity is the next
question.
\end{claim}

\subsection{Sensitivity is not fidelity}
\label{sec:methods-holdout}
\label{sec:res-generalisation}

A model that responds to the drug token has cleared a low bar, because reproducing a training
average requires responding to it too. Everything scored above came from conditions the model
trained on, so the result so far shows that the token is consulted --- not that anything
transferable was learned. Separating those two requires withholding data, and the physical design
of the atlas forces which withholding is meaningful, so the design comes first.

Tahoe applies one drug at one dose to one well containing $\approx 49$ cell lines simultaneously
--- measured: $289$ treatment samples, a mean of $48.6$ cell lines per well (median $49$, range
$43$--$50$), and $286$ of $287$ distinct (drug, dose) treatments occurring in exactly one well.
Three consequences follow, and they apply to every analysis of this atlas, not only this one. A
drug cannot be present in cell line A and absent from cell line B at the same dose, because the two
lines are physically the same well. So a $(\textrm{drug}, \textrm{cell line})$ holdout --- the
field's standard Leave-Pairs-Out split, and the split an earlier version of this work used ---
leaves the held-out condition's treated well sitting in the training set through its other
$\approx 48$ cell lines. Measured against the holdout this investigation previously used, $124$ of
the $268$ treated wells in this cache placed a held-out condition and a training condition in the
same physical sample. What that estimand measures is \emph{within-atlas matrix completion}, not
independent generalisation, and no processing of the targets changes it, because it is a property
of what was randomised.

Holding out \emph{whole wells} is the only assignment-respecting choice the design offers, and it
answers a different and narrower question: generalisation to an \emph{unseen treatment well of a
seen drug} --- a new physical preparation, not a new cellular context. Fixed-dose cross-context
transfer, present in one line and absent in another, is structurally unidentifiable in this atlas;
this is a property of Tahoe's design, not of our pipeline, and no rebuild fixes it.

The split used throughout the residual arm is therefore keyed at the well. \texttt{train} holds
every condition of the remaining wells ($4{,}999$ conditions before the reliability line,
$2{,}523$ after it). \texttt{unseen\_combo} holds $36$ treated wells out whole ($904$ conditions,
$900$ after inventory filters), and the holdout never takes a drug's last training well --- without
that guard a two-well drug could silently become an unseen \emph{drug} while sitting in the
unseen-combination arm, making the two arms measure the same thing; zero wells contribute
conditions to both sides of this split. \texttt{unseen\_drug} holds every condition of $10$
held-out drugs ($714$ conditions, $711$ after inventory filters): the model never sees the token,
so this is the \emph{control} --- naming a drug the model knows nothing about versus naming
something else can make no difference, and a non-zero effect there is an instrument fault, not a
biological finding. We do not build a Leave-Cell-Lines-Out split here;
Section~\ref{sec:limitations} argues it is the natural next axis and explains why its interest is
contingent on the transfer coefficient. Two warnings ride with the design. DrEval reports,
citing~\cite{partin}, that apparently good performance on pair-holdout splits can be reached simply
by predicting each drug's average response across the training cell lines; our held-out-well
estimand sits exactly in the regime where that warning bites, which is why the baseline ladder of
\S\ref{sec:res-lookup} is not optional --- the drug-average predictor is not a formality, it is the
hypothesis under test. And evaluation samples conditions under an explicit per-split quota, because
a uniform draw reproduces the pool's proportions and starves precisely the split the experiment
exists to test: in a first run, $250$ constructed held-out combinations yielded only $33$ scored
conditions, purely as an artefact of the sampler. The residual-frame evaluation therefore scores
$200$ train, $250$ \texttt{unseen\_combo} and $120$ \texttt{unseen\_drug} conditions --- the latter
two being quotas of the $900$ and $711$ available, so the held-out verdicts below are powered on a
subset of the conditions the split affords, and their detection limits should be read with that in
mind.

The first finding from this design is about the comparator, and it reframes every gap quoted so
far. The comparator's own score gives it away. Across the three swap strata, the scrambled model's
$\NIR$ is $0.572$, $0.497$ and $0.403$ as the named partner's true response moves from aligned
($\cos = +0.26$) through orthogonal ($0.00$) to anti-aligned ($-0.24$) with the target. A
comparator whose score is a monotone function of which drug it was handed is not a null: it is an
active treatment that tracks the partner drug. The swap partner is itself a drug the model was
mostly trained on, so telling the model a lie costs it precisely because it knows the partner. A
gap measured against such an arm decomposes as
\begin{equation}
  \gap \;=\; \underbrace{\big(\NIR_{\textrm{model}} - 0.5\big)}_{\textrm{how much naming the truth helps}}
  \;+\; \underbrace{\big(0.5 - \NIR_{\textrm{scramble}}\big)}_{\textrm{how much naming the lie hurts}},
\end{equation}
and only the first term is wanted. Against \texttt{opposite} the second term is $+0.097$; against
\texttt{orth}, whose own score sits on chance ($0.497$), it is zero. Only the neutral column of
Table~\ref{tab:generalisation} can therefore carry a magnitude claim, and the
\texttt{opposite} column is reported solely as the diagnostic that would have misled us.

\begin{table}[h]
\centering
\begin{tabular}{lrrll}
\toprule
Split & $n$ (cell lines) & model $\NIR$ & $\gap$ vs \texttt{orth} (neutral) & $\gap$ vs \texttt{opposite} \\
\midrule
\texttt{train} & 200 (35) & 0.635 [0.575, 0.695] & \ci{+0.1457}{+0.0813}{+0.2101} & \ci{+0.2627}{+0.2081}{+0.3174} \\
\textbf{\texttt{unseen\_combo}} & \textbf{250 (39)} & \textbf{0.545 [0.493, 0.596]} & \ci{+0.0449}{-0.0095}{+0.0994} & \ci{+0.1281}{+0.0747}{+0.1816} \\
\texttt{unseen\_drug} & 120 (38) & 0.512 [0.442, 0.581] & \ci{+0.0097}{-0.0694}{+0.0888} & \ci{+0.0862}{-0.0028}{+0.1753} \\
\bottomrule
\end{tabular}
\caption{Whole-well holdout on the rebuilt targets (split-before-fit, plate-scoped generic). Only
the \texttt{orth} column can carry a magnitude claim; the \texttt{opposite} column adds ``how much
naming a lie hurts'' to ``how much naming the truth helps'' and is shown as the diagnostic. The
\texttt{unseen\_drug} row is the control that proves the point: those are drugs the model has never
seen, its gap must be zero, and it is zero only in the neutral column. Model $\NIR$ intervals test
against the chance value $0.5$ with no comparator. Intervals are two-way cluster-robust over cell
lines and wells, with a Student $t$ reference (\S\ref{sec:methods-data}).}
\label{tab:generalisation}
\end{table}

The control behaves exactly as the diagnosis predicts. On \texttt{unseen\_drug} --- drugs held out
of fine-tuning, where no learned knowledge can be at work --- the gap against \texttt{opposite} is
$+0.0862$ \nolinebreak[3]$[-0.0028, +0.1753]$, nine times the neutral estimate, with an interval
that only just contains zero. The same conditions, the same generations and the same model scored
against the neutral stratum give $+0.0097$ \nolinebreak[3]$[-0.0694, +0.0888]$, sitting on zero as
a control should. The borderline interval on the \texttt{opposite} arm is itself a lesson in the
inference regime: a cluster-robust variance is asymptotic in the number of \emph{clusters} rather
than observations, this arm spans only $25$ treatment wells, and under the correct $t_{24}$
reference ($2.06$ rather than $1.96$) the interval widens by five per cent and covers zero. An
earlier draft read the normal-quantile interval as excluding zero; it does not.

Read in the neutral column, \texttt{train} clears zero at \ci{+0.1457}{+0.0813}{+0.2101} and
\texttt{unseen\_combo} does not: $+0.0449$ \nolinebreak[3]$[-0.0095, +0.0994]$, spanning zero
under well-clustering and line-clustering alike ($[-0.0050, +0.0993]$). Generalisation to an unseen
treatment well of a seen drug is therefore \emph{not established}. The precision matters: not
established is not the same as shown to be zero. The interval reaches $+0.099$, so an effect of up
to roughly $+0.1$ remains compatible with these data at $n = 250$ --- a limit of power at the
sample size the holdout affords, reported as one. (On \texttt{unseen\_drug} the corresponding
detection limit is $+0.099$.) The difficulty of the three splits is not the explanation: their
split-half replicate ceilings are $0.962$, $0.955$ and $0.956$, so the splits are equally hard and
the held-out shortfall is not a noise-floor difference.

The difference between the trained and held-out estimates is $+0.1008$
\nolinebreak[3]$[+0.0161, +0.1854]$ with cluster-robust $p = 0.0212$ (an unclustered permutation
returns $p = 0.0054$, quoted only for contrast, because ignoring the clustering this chapter
mandates everywhere else flatters the result).
% AUDITOR NOTE: the [+0.0196, +0.1820] interval is carried over from the pre-revision text; the
% fixed-claims register fixes only the point estimate and p. If the interval was computed under a
% normal rather than Student-t reference it should be recomputed before submission.
What that advantage \emph{is} requires care, because the two arms differ in two ways at once:
whether the condition was trained on, and \emph{which drugs they contain}. Restricting the
comparison to the $20$ drugs present in both arms, the advantage attenuates to $+0.0420$ with
$p = 0.4851$. The defensible statement is therefore narrow: a condition-weighted training-exposure
advantage exists in this evaluation, and its generality across drugs is unresolved --- the design
cannot separate exposure from drug identity, and we do not claim either direction of the stronger
interpretation.

Two instruments that need no comparator at all cannot be touched by any of this, and they carry the
positive content of the section. First, the slope test: regressing each generation's $\NIR$ on the
cosine between the named drug's truth and the target's truth gives \ci{+0.390}{+0.225}{+0.542} on
\texttt{train}, \ci{+0.374}{+0.214}{+0.508} on \texttt{unseen\_combo} and
\ci{+0.238}{+0.033}{+0.420} on \texttt{unseen\_drug} --- all three exclude zero. The model emits
something aligned with the named drug's signature on every split, including held-out wells, with no
degradation worth the name between train and held-out. Second, scored against chance with no
comparator, the model's own $\NIR$ clears $0.5$ on \texttt{train} ($0.635$ [$0.575$, $0.695$]) and
is not distinguishable from it on either held-out split. A positive slope alongside a chance-level
$\NIR$ is not a contradiction: it says the drug-to-signature map exists but is \emph{coarse}. The
model knows roughly what a drug does and not what it does \emph{in this well} --- and $\NIR$ asks
the second question, ranking against the $\approx 40$ other drugs measured in the same cell line.
That is the drug main effect without the drug $\times$ cell-line interaction, which
\S\ref{sec:res-transfer} measures at $44.7\%$ of the residual variance. It is also why the lookup
wins in the next section: \texttt{drug\_lookup} \emph{is} that main effect, estimated straight from
data far more precisely than the model learned it.

The directional-agreement measurement says the same thing at condition level. The cosine between
the model's prediction and its own condition's truth on trained conditions averages $+0.0511$
(sd $0.1141$, median $+0.0300$), with individual conditions spanning $[-0.1819, +0.4395]$ and
$32.0\%$ clearing $|0.1|$: weak on average and substantially heterogeneous, so a single summary
number would misrepresent it in either direction.

One caveat attaches to every held-out number in this chapter, and it is stated here where those
numbers live. The evaluation's truths are reconstructed through a compatibility path that refits
the generic on the reliability-surviving subset of training conditions rather than the full set the
checkpoint was trained against, and that filters held-out conditions by their own split-half
reproducibility --- a selection on the evaluation set's outcome. Its direction is conservative for
the headline: it retains the cleanest, easiest conditions, and the held-out result still fails to
clear zero. But the evaluated population is not the full held-out population, and the intervals
above should be read as applying to the reproducible subset.

Two instrument defects, both of which would have concealed everything in this section, are worth
recording because they are now checks rather than assumptions. A uniform condition sampler
reproduces the pool's split proportions, so $250$ constructed held-out combinations yielded only
$33$ scored ones; and a strict-inequality $\NIR$ scored the zero vector (``predict the generic
response'') at $0.000$ rather than $0.500$. Both are corrected (\S\ref{sec:methods-metrics},
\S\ref{sec:methods-holdout}), and after correction the drug-agnostic baselines land where they
must: \texttt{generic} exactly $0.500$, \texttt{control\_copy} $0.511$, \texttt{random} $0.521$.
The residual frame is \emph{demonstrably} leak-proof rather than assumed to be.

\begin{claim}
The sentence the evidence supports is narrower than the one this section originally carried:
re-encoding the target made the model respond to the combined drug-and-mechanism prompt and learn a
drug-level average response. It learned the main effect and not the interaction, and its
generalisation to unseen wells is a power-limited open question rather than a result in either
direction.
\end{claim}

\subsection{A lookup table wins}
\label{sec:methods-baselines}
\label{sec:res-lookup}

Transfer without fidelity is a real property, but it says nothing about quality: a model can apply
a learned signature in a new context and still apply it badly. The pre-registered bar was therefore
never ``beats chance'' --- residual targets are plausibly memorisable, and a predictor that simply
looks up what this drug did elsewhere needs no model at all. That predictor is the comparison that
matters, and it needs a fair construction. Each baseline below is a predictor of the drug-specific
residual, scored through an identical pipeline --- same truths, same signed-rank encoding, same
comparison set --- and two disciplines make the ladder fair. The lookups are fitted from
\textbf{training conditions only}: fitted from the full cache, a lookup scoring a held-out
condition reads residuals the model was never shown --- an oracle, not a baseline --- and the
oracle variants are reported alongside, explicitly labelled, as a bound on how much the comparison
was ever inflated rather than hidden. And \texttt{drug\_lookup\_1} selects a cell \emph{line} and
averages that line's conditions rather than selecting a single condition, because $62\%$ of drugs
are multi-dose and a single condition would confound ``different cell line'' with ``different
dose''. \texttt{moa\_lookup} is not an oracle in any split, as it uses only drugs that were in
training.

\begin{table}[h]
\centering
\begin{tabular}{lll}
\toprule
Baseline & Information used & Role \\
\midrule
\texttt{drug\_lookup}   & this drug's mean residual, all other cell lines & the bar to beat \\
\texttt{drug\_lookup\_1} & this drug in \emph{one} other cell line & removes the averaging advantage \\
\texttt{moa\_lookup}    & other drugs sharing this drug's MoA, same cell line & mechanism-leak diagnostic \\
\texttt{control\_copy}  & none (predict the control cell) & must be $\approx 0.50$ \\
\texttt{generic}        & none (predict the average response) & must be $\approx 0.50$ \\
\texttt{random}         & none & floor \\
\bottomrule
\end{tabular}
\caption{The baseline ladder. \texttt{drug\_lookup} is the pre-registered win condition: a lookup
can reproduce a memorised average, so only a model that tailors the response to the control cell's
state can beat it.}
\label{tab:baselines}
\end{table}

\begin{table}[h]
\centering
\begin{tabular}{lrl}
\toprule
Arm & $\NIR$ & note \\
\midrule
ceiling (within-condition split-half replicate) & $0.958$ & $0.962$ / $0.955$ / $0.956$ by split \\
\texttt{drug\_lookup} (training-only) & $0.941$ & the bar to beat \\
\texttt{drug\_lookup\_1} (one other line) & $0.859$ & removes the averaging advantage \\
model & $0.569$ & $0.635$ / $0.545$ / $0.512$ by split \\
\texttt{random} & $0.521$ & floor \\
\texttt{control\_copy} & $0.511$ & drug-agnostic \\
\texttt{generic} & $0.500$ & the zero vector, by construction \\
\bottomrule
\end{tabular}
\caption{The baseline ladder on the rebuilt targets, $n = 570$ scored conditions; the lookups are
fitted from training conditions only. The paired comparison on the common support ($n = 449$) is
model $-$ \texttt{drug\_lookup} $= \ci{-0.3559}{-0.3988}{-0.3130}$, and of the distance from chance
to the replicate ceiling the model covers $18.6\%$ against the lookup's $96.3\%$.}
\label{tab:lookup}
\end{table}

The pre-registered interpretation, written before the result, anticipated three branches; the middle
one fired. The model reads and transfers the drug, but recovers only a small fraction of its
signature. On the earlier build it also beat \texttt{moa\_lookup} ($0.529$, $n = 143$), so its
prompt response is not reducible to the \texttt{Mechanism:} field --- though at that level the bar
is low.

This is DrEval's finding, in transcriptomic space. DrEval reports that deep learning models barely
outperform a naive predictor of mean drug and cell-line effects~\cite{dreval}. Their
\texttt{NaiveMeanEffectsPredictor} is $\hat{y}_{ij} = \mu^c_i + \mu^d_j - \mu$; our
\texttt{drug\_lookup} is the $\beta(d)$ term of that same additive model, measured in a frame where
the $\textrm{generic}(c)$ term has already been subtracted. The finding replicates and sharpens: on
a $946$-gene transcriptomic response rather than a scalar potency, and against a proper
within-condition replicate ceiling, the additive drug main effect reaches $0.941$ of a $0.958$
ceiling while a $1$B-parameter model reaches $0.569$. Note also that the warning DrEval inherits
from~\cite{partin} --- that good pair-holdout performance is reachable by predicting each drug's
training average --- is treated here as the hypothesis rather than the caveat, and measured. It
holds.

It is tempting to read $\textrm{ceiling} - \texttt{drug\_lookup} = +0.017$ as ``the modelling
target has vanished''. That inference is not supported, for three reasons that all favour the
lookup: the ceiling is scored against a \emph{half-sample} truth while every baseline is scored
against the \emph{full-sample} truth; \texttt{drug\_lookup} averages $\approx 38$ conditions against
the ceiling's single noisy half; and $\NIR$ near $0.96$ is compressive, collapsing a wide range of
cosines into a few thousandths.

\begin{claim}
The pre-registered middle branch fired: the model reads and transfers the drug, but recovers only a
small fraction of its signature. The honest summary is the paired gap on common support, $-0.3559$
\nolinebreak[3]$[-0.3988, -0.3130]$ --- $18.6\%$ of the chance-to-ceiling distance covered, against
the lookup's $96.3\%$ --- and it is DrEval's finding replicated in transcriptomic space: the
additive drug main effect, measured against a proper within-condition ceiling, reaches $0.941$ of
$0.958$ where a $1$B-parameter model reaches $0.569$.
\end{claim}

\subsection{What is closed and open for unseen drugs}
\label{sec:res-scope}

An unseen drug can only be predicted through a property it \emph{shares} with drugs seen in
training, so the question is answerable by enumerating the channels and measuring each in the same
frame. For a channel $X$, a drug's residual is predicted as the mean residual of the \emph{other}
drugs in the same cell line that share property $X$ with it --- the construction of
\texttt{moa\_lookup}, generalised. Because it uses only drugs that were in training, it is a fair
contest in every split, unlike a same-drug lookup, which on an unseen drug is an oracle. The gate's
estimand requires the channel arm and its null to differ in exactly one respect, and two
disciplines enforce it. Every channel is paired with a \textbf{plate-matched random null}: the same
\emph{number} of partner drugs, drawn at random from the same cell line, \emph{and} the same split
of partners across the target's own plate and other plates, because mechanism- and target-related
drugs are co-plated more often than chance --- a null matched on count alone would differ from the
channel arm in plate membership as well as in the channel, which is two differences where the
estimand allows one. A \textbf{co-plating diagnostic is printed before any verdict}, so the size of
the correction is a number rather than an argument, and where a plate-matched null spans too few
cell lines the channel is reported \textsc{untestable} rather than closed --- ``we could not build
the control'' must never be written as ``closed''. Coverage is likewise reported before any
verdict, because a channel annotated on few drugs cannot be closed by a null result: it was never
tested.

\begin{table}[h]
\centering
\begin{tabular}{lrrrrl}
\toprule
Channel & Coverage & $n$ & vs plate-matched null & (vs count-matched) & Verdict \\
\midrule
\textbf{protein target} & 54\% & 877 & \ci{+0.0973}{+0.0741}{+0.1183} & $+0.1269$ & \textbf{live} \\
\textbf{mechanism (MoA)} & 32\% & 1033 & \ci{+0.0749}{+0.0538}{+0.0948} & $+0.1067$ & \textbf{live} \\
chemical structure & 100\% & 4074 & $+0.0162$ [${+0.0051}$, ${+0.0264}$] & $+0.0206$ & closed \\
\bottomrule
\end{tabular}
% AUDITOR NOTE: the coverage and n columns are carried over from the count-matched run; the
% plate-matched re-run's register does not restate them. They should be near-identical (same
% annotation support) but ought to be re-quoted from the corrected artifact before submission.
\caption{Coverage is stated first: a channel annotated on few drugs cannot be \emph{closed} by a
null result, because it was never tested. The plate-matched null is the estimand-correct control;
the count-matched column shows how much of each raw gap was co-plating. Chemistry reproduces the
structure gate in this frame; both knowledge channels clear their corrected nulls.}
\label{tab:channels}
\end{table}

The size of that correction is itself informative, so it is reported rather than argued away.
Mechanism partners are co-plated with the target $0.420$ of the time against $0.340$ for a random
draw (protein target $0.348$ against $0.323$; chemical similarity $0.344$ against $0.345$ --- none,
as expected for a channel uncorrelated with plating). Under a third construction that draws all
partners from \emph{other} plates, the knowledge channels' advantages fall to $+0.0643$ and
$+0.0344$ and chemistry's interval spans zero ($+0.0056$); and in that construction the channel
arms' absolute $\NIR$ sits at $0.360$ and $0.296$, \emph{below} chance --- a direct demonstration of
how strongly the cell-line-scoped residual frame is dominated by plate. The verdicts above survive
the correction; the magnitudes shrank by a fifth to a third, which is what the co-plating fractions
predict. The bound this establishes is on \emph{similarity transfer} through each channel, not on
any model: a model could in principle combine channels or learn a non-smooth structure-to-response
map that a nearest-neighbour average cannot express, though at a few hundred drugs the lookup is
effectively the bound.

This overturns a pre-registration made in this project. Two design documents concluded that
unseen-drug generalisation was out of reach, and both inferred it from the structure gate --- that
is, from Morgan fingerprints and a molecular language model. But the drug-knowledge specification
written earlier in the same project ranks raw structure the \emph{lowest}-value channel and protein
targets ``the key feature'', on the reasoning that transcriptional response is driven by which
protein a compound hits, something structure encodes only indirectly. The pre-registration
generalised from the weakest channel to the strongest, and it was wrong. Tahoe's own metadata
carries protein-target annotations for $264$ drugs --- $280$ distinct symbols, $550$ edges ---
which no analysis in this project had read.

\begin{claim}
The correct statement is narrower and more useful than the one we started with: chemical structure
does not predict response, measured twice in two frames; protein target and mechanism do.
Unseen-drug generalisation is reachable through drug-side \emph{knowledge}, and not through
chemistry.
\end{claim}

Three caveats bound it. Coverage is $54\%$ and $32\%$, so this is measured
on the annotated subset. The channel gaps are severalfold below the advantage a same-drug lookup
enjoys for a \emph{seen} drug ($+0.44$ above chance), so the reachable ceiling in this regime is
much lower. And target and MoA overlap, so whether target adds anything \emph{over} MoA is a
further question this measurement does not settle --- settling it would need a
leave-one-mechanism-out split rather than the leave-one-drug-out split used here, since a held-out
drug's mechanism class is generally still represented in training.

\label{sec:res-field}
The natural next step writes itself: append the protein targets to the prompt, retrain, and see
whether unseen-drug performance rises. That arm assumes the bottleneck is \emph{information
availability}. Everything in this chapter points the other way, but the direct evidence is weaker
than earlier drafts claimed, and it is worth saying exactly how weak it is. The scramble has until
now replaced the drug name \emph{and} the mechanism together, so a non-null gap establishes
sensitivity to the combined prompt without attributing it to either span. A \emph{field
decomposition} costs one extra pair of generation arms and swaps exactly one span of the
instruction line to the opposite-signature partner, leaving every other byte identical: the drug
name with the mechanism kept, the mechanism with the name kept, and both together. One guard is
essential rather than tidy: if the swap partner happens to share the original's mechanism, then
swapping only the mechanism produces a prompt \emph{identical} to the model's own, and the arm
scores a gap of exactly zero \emph{by construction} --- read naively, indistinguishable from ``the
model ignores the mechanism'', the very claim under test --- so such conditions are dropped rather
than scored, which is why the mechanism arm carries fewer conditions than the others.

On the earlier target build, the decomposed arms suggested the response is carried almost entirely
by the name span: swapping only the drug name gave $\gap = \ci{+0.0809}{+0.0592}{+0.1010}$
($n = 570$), swapping only the mechanism gave $+0.0091$ [$-0.0174$, $+0.0383$] ($n = 325$), and
both together $+0.1024$ [$+0.0827$, $+0.1236$]. Those arms were scored against the
\texttt{opposite}-signature comparator that \S\ref{sec:res-generalisation} shows inflates every
gap read against it, and they were not re-run under the corrected evaluation. The claim that the
model ``reads the name and not the mechanism'' is therefore withdrawn pending that re-run; what the
current run establishes is response to the combined drug-and-mechanism prompt, with the mechanism
field's contribution unresolved.

The practical consequence is narrower than the one earlier drafts drew, and more honest. The
target-prompt arm is not retired by measurement --- it is untested. What argues for building it
carefully rather than eagerly is the shape of the failure documented above: the model's deficit is
not confined to missing information about unseen drugs, because on \emph{seen} drugs it recovers
only a fifth of the chance-to-ceiling distance against a lookup that knows nothing about contexts
at all. And what argues for building it at all is the one thing a lookup structurally cannot do:
handle a drug it has never seen, in the one regime where this section shows live channels exist.
The model's opportunity is not the interaction --- \S\ref{sec:res-transfer} found $44\%$ of the
residual variance there but no learnable structure in it --- it is the unseen drug. Whether the
model can be made to take that opportunity is the question the thesis leaves open.

\subsection{The investigation in one table}
\label{sec:res-summary}

The results are distributed across two measurement frames, and a reader assembling them has to keep
track of which numbers may be compared with which. The two frames are \emph{not} interchangeable
and are separated deliberately. The expression frame scores a predicted cell profile against other
drugs on the same plate; the residual frame scores a predicted drug-specific signature against
other drugs in the same cell line, after the generic program has been subtracted. Both have chance
at $0.50$ by construction, and that coincidence is the trap: the residual frame's ceiling is
$0.958$ where the expression frame's is $0.576$, so a model scoring $0.569$ in one and $0.498$ in
the other has not improved --- it is being measured against a different question.

\begin{table}[h]
\centering
\begin{tabular}{llrl}
\toprule
Frame & Arm & $\NIR$ & What it knows \\
\midrule
\multirow{5}{*}{\shortstack[l]{expression\\(within plate,\\tier 2, $n=606$)}}
 & noise ceiling & 0.576 & a real replicate of this condition \\
 & control-copy & 0.504 & nothing --- it reproduces the untreated cell \\
 & linear map & 0.500 & the control cell, no drug \\
 & \textbf{model} & \textbf{0.498} & control cell $+$ drug \\
 & global mean & 0.180 & nothing \\
\midrule
\multirow{7}{*}{\shortstack[l]{residual\\(cell-line sets,\\$n=570$)}}
 & noise ceiling & 0.958 & a real split-half replicate of this condition \\
 & drug lookup (training-only) & 0.941 & this drug's signature from other cell lines \\
 & drug lookup, one line & 0.859 & the same, from a single other cell line \\
 & \textbf{model} & \textbf{0.569} & control cell $+$ drug \\
 & scramble, \texttt{orth} partner & 0.497 & the neutral comparator: chance by construction \\
 & scramble, \texttt{opposite} partner & 0.403 & an active treatment, not a null \\
 & generic response & 0.500 & nothing (the zero vector, by construction) \\
\bottomrule
\end{tabular}
\caption{Every arm on the calibrated metric, by frame. Chance is $0.50$ in both. The residual-frame
figures are the rebuilt-target run of \S\ref{sec:res-generalisation} and \S\ref{sec:res-lookup};
the held-out subset carries the outcome-selection caveat stated there.}
\label{tab:allarms}
\end{table}

Read down the residual block and the investigation's shape is visible in one object. A replicate of
the experiment reaches $0.958$; a lookup table reaches $0.941$; the model reaches $0.569$; the same
model told a response-orthogonal drug sits at $0.497$, the neutral zero point, and told an
anti-aligned drug falls to $0.403$. The distance between the model and its neutral scramble is the
drug effect this work recovered; the distance between the model and the lookup is what it did not.
